% Task C08 solutions
% Based on chapters/ch08.tex from local main b4628dd90903b38fb9c71bf61921e79513bb5295.
\chapter*{第8章：微分学的应用习题解答}
\addcontentsline{toc}{chapter}{第8章：微分学的应用习题解答}
\section*{8.1.3 练习题}

\paragraph{题 1}
\begin{xsolution}
逐题检查 L'Hospital 法则的适用条件及求导后的表达式。
\begin{enumerate}[label=(\arabic*)]
  \item 当 $x\to1$ 时原式趋于 $1/2$，既不是 $0/0$ 型也不是 $\infty/\infty$ 型，因此不能使用 L'Hospital 法则。

  \item 原式虽为 $\infty/\infty$ 型，但求导后得到
  \[
    \frac{1+\cos x}{1-\cos x}.
  \]
  该比值在 $x\to+\infty$ 时没有极限，而且分母导数 $1-\cos x$ 在任意远处仍有零点，故 L'Hospital 法则不能给出所求极限。直接同除以 $x$ 即得
  \[
    \frac{1+\sin x/x}{1-\sin x/x}\to1.
  \]

  \item 求导后得到
  \[
    \frac{\ee^x-\ee^{-x}}{\ee^x+\ee^{-x}},
  \]
  恰为原比值的倒数；若继续求导又回到原式，因此单靠反复使用 L'Hospital 法则只会循环。直接同除以 $\ee^x$ 得
  \[
    \frac{1+\ee^{-2x}}{1-\ee^{-2x}}\to1.
  \]

  \item 求导后得到
  \[
    \frac{1}{x/\sqrt{1+x^2}}=\frac{\sqrt{1+x^2}}x,
  \]
  仍是原比值的倒数，因而同样形成循环。直接同除以 $x$：
  \[
    \frac{x}{\sqrt{1+x^2}}=\frac1{\sqrt{1+x^{-2}}}\to1.
  \]
\end{enumerate}
\end{xsolution}

\paragraph{题 2}
\begin{xsolution}
\begin{enumerate}[label=(\arabic*)]
  \item 设 $a,b>0$ 且 $a\ne b$。当 $x\to0$ 时
  \[
    a^{x^2}-b^{x^2}=x^2(\ln a-\ln b)+\littleo(x^2),
  \]
  而
  \[
    a^x-b^x=x(\ln a-\ln b)+\littleo(x).
  \]
  因而
  \[
    \lim_{x\to0}\frac{a^{x^2}-b^{x^2}}{(a^x-b^x)^2}
    =\frac1{\ln a-\ln b}=\frac1{\ln(a/b)}.
  \]
  若 $a=b$，原分母恒为零，原式没有通常意义。

  \item 令 $t=1/x$。有
  \[
    (1-t+\tfrac12t^2)\ee^t=1+\frac16t^3+\bigO(t^4),
  \]
  故
  \[
    \left(x^3-x^2+\frac x2\right)\ee^{1/x}
    =x^3+\frac16+\bigO(x^{-1}).
  \]
  又
  \[
    \sqrt{1+x^6}=x^3\sqrt{1+x^{-6}}=x^3+\bigO(x^{-3}).
  \]
  所以极限为 $1/6$。

  \item 取对数：
  \[
    \ln\left(\ln\frac1x\right)^x=x\ln\ln\frac1x.
  \]
  令 $t=\ln(1/x)\to+\infty$，则右端为 $\ee^{-t}\ln t\to0$，故极限为 $1$。

  \item 因 $x>0$ 时 $\frac\pi2-\arctan x=\arctan(1/x)\sim1/x$，故
  \[
    \ln\left(\frac\pi2-\arctan x\right)^{1/\ln x}
    =\frac{\ln(\frac\pi2-\arctan x)}{\ln x}\to-1.
  \]
  因此极限为 $\ee^{-1}$。
\end{enumerate}
\end{xsolution}

\paragraph{题 3}
\begin{xsolution}
逐项比较 Maclaurin 展开式的低阶系数即可。
\begin{enumerate}[label=(\arabic*)]
  \item
  \[
    (a+b\cos x)\sin x
    =(a+b)x+\left(-\frac{a+b}{6}-\frac b2\right)x^3
    +\left(\frac{a+b}{120}+\frac b8\right)x^5+\bigO(x^7).
  \]
  令一次、三次项在 $x-(a+b\cos x)\sin x$ 中消失，得
  \[
    a+b=1,\qquad \frac16+\frac b2=0,
  \]
  即
  \[
    a=\frac43,\qquad b=-\frac13.
  \]
  此时 $f(x)=\frac1{30}x^5+\bigO(x^7)$，故最高可做到五阶。

  \item
  \[
    \frac{1+ax}{1+bx}
    =1+(a-b)x-b(a-b)x^2+b^2(a-b)x^3+\bigO(x^4).
  \]
  与 $\ee^x$ 比较，令一次、二次项相同：
  \[
    a-b=1,\qquad -b=\frac12.
  \]
  因而
  \[
    a=\frac12,\qquad b=-\frac12.
  \]
  此时三次项之差为 $1/6-1/4=-1/12$，故 $f(x)=-x^3/12+\bigO(x^4)$。

  \item
  \[
    \frac{1+ax^2}{x+bx^3}
    =\frac1x+(a-b)x-b(a-b)x^3+b^2(a-b)x^5+\bigO(x^7),
  \]
  而
  \[
    \cot x=\frac1x-\frac x3-\frac{x^3}{45}-\frac{2x^5}{945}+\bigO(x^7).
  \]
  令 $a-b=-1/3$、$-b(a-b)=-1/45$，得到
  \[
    b=-\frac1{15},\qquad a=-\frac25.
  \]
  此时五次项之差为 $-1/1575$，所以最高为五阶无穷小。

  \item
  \[
    \frac{1+ax^2}{1+bx^2}
    =1+(a-b)x^2-b(a-b)x^4+b^2(a-b)x^6+\bigO(x^8).
  \]
  与
  \[
    \cos x=1-\frac{x^2}{2}+\frac{x^4}{24}-\frac{x^6}{720}+\bigO(x^8)
  \]
  比较，得
  \[
    a-b=-\frac12,\qquad \frac b2=\frac1{24},
  \]
  即
  \[
    b=\frac1{12},\qquad a=-\frac5{12}.
  \]
  此时 $f(x)=x^6/480+\bigO(x^8)$，故最高为六阶。
\end{enumerate}
\end{xsolution}

\paragraph{题 4}
\begin{xsolution}
\begin{enumerate}[label=(\arabic*)]
  \item 由 $\sin x=x-x^3/6+\bigO(x^5)$，
  \[
    \frac1x-\frac1{\sin x}
    =\frac{\sin x-x}{x\sin x}\sim-\frac x6\to0.
  \]

  \item
  \[
    \ee^{x^3}-1-x^3=\frac12x^6+\bigO(x^9),\qquad
    \sin^6 2x=64x^6+\littleo(x^6),
  \]
  故极限为 $1/128$。

  \item
  \[
    n\sin\frac1n=1-\frac1{6n^2}+\bigO(n^{-4}),
  \]
  因而
  \[
    n^2\ln\left(n\sin\frac1n\right)
    =n^2\left(-\frac1{6n^2}+\bigO(n^{-4})\right)\to-\frac16.
  \]

  \item
  \[
    \sqrt{n^2+2}\,\pi
    =n\pi+\frac\pi n+\bigO(n^{-3}).
  \]
  因此
  \[
    (-1)^n n\sin(\sqrt{n^2+2}\,\pi)
    =n\sin\left(\frac\pi n+\bigO(n^{-3})\right)\to\pi.
  \]

  \item 展开到七次项可得
  \[
    \sin(\tan x)-\tan(\sin x)=-\frac1{30}x^7+\bigO(x^9),
  \]
  故极限为 $-1/30$。

  \item 展开到六次项可得
  \[
    x\sin(\sin x)-\sin^2x=\frac1{18}x^6+\bigO(x^8),
  \]
  故极限为 $1/18$。
\end{enumerate}
\end{xsolution}

\paragraph{题 5}
\begin{xsolution}
由 $f''(a)$ 存在，
\[
  f(a+h)=f(a)+f'(a)h+\frac12f''(a)h^2+\littleo(h^2),
\]
\[
  f(a+2h)=f(a)+2f'(a)h+2f''(a)h^2+\littleo(h^2).
\]
相减后
\[
  f(a+2h)-2f(a+h)+f(a)=f''(a)h^2+\littleo(h^2),
\]
除以 $h^2$ 并令 $h\to0$ 即得结论。
\end{xsolution}

\paragraph{题 6}
\begin{xsolution}
记 $t=x-x_0$，$A=f'(x_0)\ne0$，$B=f''(x_0)/2$。二阶 Taylor 公式给出
\[
  f(x)-f(x_0)=At+Bt^2+\littleo(t^2).
\]
于是
\begin{align*}
  &\frac1{f(x)-f(x_0)}-\frac1{f'(x_0)(x-x_0)}\\
  &\quad=\frac{At-[At+Bt^2+\littleo(t^2)]}
  {At[At+Bt^2+\littleo(t^2)]}
  \to-\frac{B}{A^2}.
\end{align*}
故所求极限为
\[
  -\frac{f''(x_0)}{2[f'(x_0)]^2}.
\]
\end{xsolution}

\paragraph{题 7}
\begin{xsolution}
切线方程为
\[
  Y-f(x)=f'(x)(X-x),
\]
故其 $x$ 轴截距
\[
  u=x-\frac{f(x)}{f'(x)}.
\]
记 $c=f''(0)>0$。当 $x\to0+$ 时
\[
  f(x)=\frac c2x^2+\littleo(x^2),\qquad
  f'(x)=cx+\littleo(x),
\]
所以
\[
  u=\frac x2+\littleo(x),\qquad \frac ux\to\frac12.
\]
又 $u\to0+$，故
\[
  \frac{f(u)}{f(x)}\sim\frac{u^2}{x^2}.
\]
从而
\[
  \frac{x f(u)}{u f(x)}\sim\frac ux\to\frac12.
\]
\end{xsolution}

\paragraph{题 8}
\begin{xsolution}
记
\[
  a_n=\left(1+\frac1n\right)^n.
\]
由
\[
  n\ln\left(1+\frac1n\right)
  =1-\frac1{2n}+\frac1{3n^2}+\bigO(n^{-3})
\]
得到
\[
  a_n=\ee\left(1-\frac1{2n}+\frac{11}{24n^2}+\bigO(n^{-3})\right).
\]
因此
\[
  a_{n+1}-a_n=\frac{\ee}{2n^2}+\bigO(n^{-3}).
\]
于是
\[
  n^x(a_{n+1}-a_n)
  =\frac\ee2 n^{x-2}(1+\littleo(1)).
\]
按通常实值函数的定义，极限有限的 $x$ 恰为 $x\le2$，且
\[
  f(x)=
  \begin{cases}
    0,&x<2,\\
    \ee/2,&x=2.
  \end{cases}
\]
故定义域为 $(-\infty,2]$，值域为 $\{0,\ee/2\}$。若允许扩充实数值，则 $x>2$ 时极限为 $+\infty$。
\end{xsolution}

\paragraph{题 9}
\begin{xsolution}
因 $0<\ln(1+t)<t$（$t>0$），有 $0<x_{n+1}<x_n$，故 $x_n\to L\ge0$。令极限进入递推式得 $L=\ln(1+L)$，所以 $L=0$。

考察 $1/x_n$，有
\begin{align*}
  \frac1{x_{n+1}}-\frac1{x_n}
  &=\frac{x_n-\ln(1+x_n)}{x_n\ln(1+x_n)}\\
  &\to\frac12,
\end{align*}
其中用了
\[
  \ln(1+x)=x-\frac{x^2}{2}+\littleo(x^2).
\]
由 Stolz 定理，
\[
  \lim_{n\to\infty}\frac{1/x_n}{n}=\frac12,
\]
即
\[
  \boxed{\lim_{n\to\infty}nx_n=2}.
\]
\end{xsolution}

\paragraph{题 10}
\begin{xsolution}
由定义
\[
  x_{n+1}=x_n+\ln(a-x_n).
\]
令 $y_n=a-x_n$，则
\[
  y_{n+1}=y_n-\ln y_n,
  \qquad y_0=a-\ln a.
\]
由基本不等式 $\ln a\le a-1$ 得 $y_0\ge1$。若 $y_n>1$，则
\[
  1<y_n-\ln y_n<y_n,
\]
其中左边由函数 $u-\ln u$ 在 $u\ge1$ 上递增且在 $u=1$ 取值 1 得到。因此 $\{y_n\}$ 单调减少且以下界 1 有界；若其极限为 $L$，则
\[
  L=L-\ln L,
\]
故 $L=1$。当 $a=1$ 时从一开始即 $y_n=1$。综上
\[
  \boxed{\lim_{n\to\infty}x_n=a-1}.
\]
\end{xsolution}

\section*{8.2.2 练习题}

\paragraph{题 1}
\begin{xsolution}
两个方向都不成立。

单调可微函数的导函数不必单调。例如
\[
  f(x)=x+\frac1{10}\sin x
\]
满足 $f'(x)=1+\frac1{10}\cos x>0$，故 $f$ 在 $\RR$ 上严格增加，但 $f'$ 不是单调函数。

反过来，导函数单调也不能推出原函数在整个区间单调。例如 $f(x)=x^2$ 的导函数 $f'(x)=2x$ 严格增加，但 $f$ 在 $\RR$ 上先减后增。
\end{xsolution}

\paragraph{题 2}
\begin{xsolution}
令
\[
  F(x)=x\ln\left(1+\frac1x\right).
\]
在 $(-\infty,-1)\cup(0,+\infty)$ 上
\[
  F'(x)=\ln\left(1+\frac1x\right)-\frac1{x+1}.
\]
令 $t=1/x$，则 $t>-1$ 且 $t\ne0$。由例题 8.5.1 的不等式
\[
  \ln(1+t)>\frac{t}{1+t}
\]
得到 $F'(x)>0$。由于原函数为 $\ee^{F(x)}$，故在两个区间上均严格单调增加。
\end{xsolution}

\paragraph{题 3}
\begin{xsolution}
取 $0<x<y$。因为 $f'$ 严格单调增加，由 Lagrange 中值定理，存在
\[
  \xi_1\in(0,x),\qquad \xi_2\in(x,y)
\]
使
\[
  \frac{f(x)-f(0)}x=f'(\xi_1),
  \qquad
  \frac{f(y)-f(x)}{y-x}=f'(\xi_2).
\]
于是
\[
  \frac{f(x)}x<\frac{f(y)-f(x)}{y-x}.
\]
交叉相乘整理即得
\[
  \frac{f(x)}x<\frac{f(y)}y.
\]
故 $f(x)/x$ 在 $(0,+\infty)$ 上严格单调增加。
\end{xsolution}

\paragraph{题 4}
\begin{xsolution}
设
\[
  R(x)=\frac{f(x)}{g(x)},\qquad r(x)=\frac{f'(x)}{g'(x)}.
\]
由 $g'>0$、$g(0)=0$ 知 $g(x)>0$（$0<x<a$）。对区间 $[0,x]$ 用 Cauchy 中值定理，存在 $\xi\in(0,x)$ 使
\[
  R(x)=\frac{f(x)-f(0)}{g(x)-g(0)}=r(\xi).
\]
因 $r$ 单调增加，$r(\xi)\le r(x)$。于是
\[
  R'(x)=\frac{g'(x)}{g(x)}[r(x)-R(x)]\ge0.
\]
故 $f/g$ 单调增加。
\end{xsolution}

\paragraph{题 5}
\begin{xsolution}
由 $f''(x)\le0$，$f'$ 在 $[a,+\infty)$ 上单调减少，因此
\[
  f'(x)\le f'(a)<0,\qquad x>a.
\]
所以 $f$ 严格单调减少，零点至多一个。又由 Lagrange 中值定理
\[
  f(x)\le f(a)+f'(a)(x-a)\to-\infty,
\]
而 $f(a)>0$，故连续性保证至少存在一个零点。于是零点恰有一个。
\end{xsolution}

\paragraph{题 6}
\begin{xsolution}
由于 $f'(x)>k>0$，$f$ 严格单调增加，故零点至多一个。令
\[
  b=a-\frac{f(a)}k>a.
\]
由中值定理，对于某个 $\xi\in(a,b)$，
\[
  f(b)-f(a)=f'(\xi)(b-a)>k\left(-\frac{f(a)}k\right)=-f(a),
\]
所以 $f(b)>0$。结合 $f(a)<0$，零点存在于 $(a,b)$，且因严格单调而唯一。
\end{xsolution}

\paragraph{题 7}
\begin{xsolution}
将方程写成
\[
  H(x)=(1+x+x^2/2)\ee^{-x}=a.
\]
因为二次多项式 $1+x+x^2/2=[(x+1)^2+1]/2>0$，且
\[
  H'(x)=-\frac{x^2}{2}\ee^{-x}\le0,
\]
并且除 $x=0$ 外严格小于零，所以 $H$ 在 $\RR$ 上严格减少。又
\[
  \lim_{x\to-\infty}H(x)=+\infty,
  \qquad
  \lim_{x\to+\infty}H(x)=0.
\]
因此对每个 $a>0$，方程 $H(x)=a$ 恰有一个实根。
\end{xsolution}

\paragraph{题 8}
\begin{xsolution}
记
\[
  c=\frac2\pi-1<0,
  \qquad
  F(x)=c\ln x-\ln2+\ln(1+x^2).
\]
则
\[
  F'(x)=\frac{c+(c+2)x^2}{x(1+x^2)}.
\]
在 $(0,1)$ 中唯一临界点为
\[
  x_0=\sqrt{\frac{\pi-2}{\pi+2}}.
\]
所以 $F$ 在 $(0,x_0)$ 上严格减少，在 $(x_0,1)$ 上严格增加。又
\[
  \lim_{x\to0+}F(x)=+\infty,
  \qquad F(1)=0.
\]
因此 $F(x)<0$ 对 $x_0<x<1$ 成立，从而 $F(x_0)<0$。于是下降支 $(0,x_0)$ 上恰穿过零点一次，而上升支在开区间 $(x_0,1)$ 内没有零点。故原方程在 $(0,1)$ 内恰有一个实根。
\end{xsolution}

\section*{8.3 思考题}

\paragraph{思考题（8.3 基本事实第 7 条）}
\begin{xsolution}
令
\[
  f(x)=
  \begin{cases}
    \ee^{-1/x^2},&x>0,\\
    0,&x=0,\\
    -\ee^{-1/x^2},&x<0.
  \end{cases}
\]
这是 $C^\infty$ 函数，并且 $f^{(n)}(0)=0$ 对每个 $n\ge0$ 都成立；这是经典的平坦函数性质，可由反复求导后每一项均为多项式乘 $\ee^{-1/x^2}$，再利用指数衰减压倒任意幂得到。

但 $x<0$ 时 $f(x)<0=f(0)$，$x>0$ 时 $f(x)>0=f(0)$，所以 $x=0$ 既不是极大值点也不是极小值点。这给出了书中所要求的“所有阶导数均为零但不取极值”的例子。
\end{xsolution}

\section*{8.3.2 练习题}

\paragraph{题 1}
\begin{xsolution}
先设 $x_0$ 是唯一极值点且为极小值点。若它不是全局最小值点，则存在 $y\in I$ 使 $f(y)<f(x_0)$。不妨设 $y>x_0$。由局部极小性，在 $x_0$ 右侧充分小的邻域内 $f(x)\ge f(x_0)$。若该邻域内恒等于 $f(x_0)$，则其中每一点都是极值点，矛盾；故可取 $z\in(x_0,y)$ 使 $f(z)>f(x_0)$。连续函数在 $[x_0,y]$ 上取得最大值，而该最大值既不能在 $x_0$ 取得，也不能在 $y$ 取得，故在内部出现另一个极大值点，矛盾。因此 $x_0$ 是最小值点。极大值情形同理。

若 $x_0$ 是极小值点而另有 $y\ne x_0$ 也取得同一最小值，则在两点之间若函数恒为该值，就有无穷多个极值点；若不恒定，则区间内部有一点取得更大的最大值，从而产生另一个极值点。故 $x_0$ 还是唯一最小值点。极大值情形同理。
\end{xsolution}

\paragraph{题 2}
\begin{xsolution}
设
\[
  F(x)=x^3+px+q.
\]
若有三个不同实根，则必须 $p<0$，否则 $F'(x)=3x^2+p\ge0$，函数单调而不可能有三个不同实根。令
\[
  s=\sqrt{-\frac p3}>0.
\]
此时 $-s$ 为极大值点，$s$ 为极小值点，且
\[
  F(-s)=q+2s^3,\qquad F(s)=q-2s^3.
\]
三条不同实根出现的充分必要条件正是
\[
  F(-s)>0,\qquad F(s)<0,
\]
即
\[
  |q|<2s^3.
\]
化为 $p,q$ 的条件为
\[
  \boxed{\,4p^3+27q^2<0\,},
\]
这本身已包含 $p<0$。
\end{xsolution}

\paragraph{题 3}
\begin{xsolution}
令
\[
  F(x)=\frac1{x^2}+px+q,\qquad x\ne0.
\]
若 $p=0$，左右两个半轴上至多各有一个正的 $x^2$ 值，不可能有三个不同实根，故必须 $p\ne0$。

有
\[
  F'(x)=p-\frac2{x^3}.
\]
若 $p>0$，则在 $(-\infty,0)$ 上 $F$ 严格增加，并从 $-\infty$ 变到 $+\infty$，恰有一个负根；在 $(0,+\infty)$ 上只有一个临界点
\[
  x_0=\left(\frac2p\right)^{1/3},
\]
且为极小值点。要在正半轴另有两个不同根，充要条件是
\[
  F(x_0)=q+3\left(\frac p2\right)^{2/3}<0.
\]
$p<0$ 时情形关于原点对称：正半轴恰有一个根，负半轴有两个根的充要条件仍为
\[
  q+3\left(\frac{|p|}{2}\right)^{2/3}<0.
\]
因此总的充分必要条件为
\[
  \boxed{p\ne0,\qquad q<-3\left(\frac{|p|}{2}\right)^{2/3}}.
\]
\end{xsolution}

\paragraph{题 4}
\begin{xsolution}
先设 $\lambda\ne0$。由算术平均值－几何平均值不等式，
\[
  a^2\ee^{\lambda x}+b^2\ee^{-\lambda x}
  \ge2ab,
\]
且等号当且仅当
\[
  a^2\ee^{\lambda x}=b^2\ee^{-\lambda x},
  \qquad
  x=\frac1\lambda\ln\frac ba.
\]
故极小值为
\[
  \boxed{2ab},
\]
与 $\lambda$ 无关。

若允许退化参数 $\lambda=0$，则函数恒为 $a^2+b^2$；因此原题“与 $\lambda$ 无关”的通常理解应当排除 $\lambda=0$（或另把这个退化情形列出）。
\end{xsolution}

\paragraph{题 5}
\begin{xsolution}
若
\[
  f(x)=\frac{ax+b}{cx+d}
\]
不是常值函数，则 $ad-bc\ne0$。在定义域的每个连通区间上
\[
  f'(x)=\frac{ad-bc}{(cx+d)^2}
\]
恒不为零且符号不变，所以 $f$ 在每个定义区间上严格单调，不存在内点极值。
\end{xsolution}

\paragraph{题 6}
\begin{xsolution}
\begin{enumerate}[label=(\arabic*)]
  \item 比较对数即可。令 $\phi(x)=\ln x/x$。因
  \[
    \phi'(x)=\frac{1-\ln x}{x^2},
  \]
  $\phi$ 在 $[\ee,+\infty)$ 上严格减少。由于 $\pi>\ee$，
  \[
    \frac{\ln\pi}{\pi}<\frac1\ee,
  \]
  即 $\ee\ln\pi<\pi$，故
  \[
    \boxed{\pi^{\ee}<\ee^\pi}.
  \]

  \item 两边取对数并约去正因子 $\sqrt{n(n+1)}/2$，问题等价于比较
  \[
    \psi(n)=\frac{\ln n}{\sqrt n},
    \qquad
    \psi(n+1)=\frac{\ln(n+1)}{\sqrt{n+1}}.
  \]
  有
  \[
    \psi'(x)=\frac{1-\frac12\ln x}{x^{3/2}},
  \]
  故 $\psi$ 在 $(0,\ee^2)$ 上增加，在 $(\ee^2,+\infty)$ 上减少。于是 $1\le n\le6$ 时
  \[
    (\sqrt n)^{\sqrt{n+1}}<(\sqrt{n+1})^{\sqrt n};
  \]
  $n\ge8$ 时不等号反向。对唯一跨过极大点的整数 $n=7$ 直接比较
  \[
    \frac{\ln7}{\sqrt7}>\frac{\ln8}{\sqrt8}
  \]
  （两边数值约为 $0.73545$ 与 $0.73519$），所以 $n=7$ 时也为反向不等号。综上，$n\le6$ 时右边较大，$n\ge7$ 时左边较大。
\end{enumerate}
\end{xsolution}

\paragraph{题 7}
\begin{xsolution}
\begin{enumerate}[label=(\arabic*)]
  \item 设三角形面积为 $K$，半周长为 $s$。令
  \[
    x=s-a,\quad y=s-b,\quad z=s-c,
  \]
  则 $x+y+z=s$，Heron 公式给出 $K^2=sxyz$。由
  \[
    xyz\le\left(\frac s3\right)^3
  \]
  得
  \[
    K^2\le\frac{s^4}{27}.
  \]
  固定 $K$ 时，$s$ 因而有确定的下界，且等号当且仅当 $x=y=z$，即 $a=b=c$。故周长最小的是等边三角形。

  \item 固定周长即固定 $s$。同一公式
  \[
    K^2=sxyz\le s\left(\frac s3\right)^3
  \]
  表明面积最大当且仅当 $x=y=z$，故仍为等边三角形。

  \item 取第一象限顶点 $(x,y)$，矩形的长、宽为 $2x,2y$，且
  \[
    \frac{x^2}{a^2}+\frac{y^2}{b^2}=1.
  \]
  令 $u=x/a$，$v=y/b$，则 $u^2+v^2=1$。面积
  \[
    4xy=4abuv\le2ab,
  \]
  等号在 $u=v=1/\sqrt2$ 时成立。因此最大面积矩形的长、宽为
  \[
    \boxed{\sqrt2\,a,\qquad \sqrt2\,b}.
  \]

  \item 周长为 $4(x+y)$。由 Cauchy 不等式
  \[
    x+y=au+bv\le\sqrt{a^2+b^2}\sqrt{u^2+v^2}=\sqrt{a^2+b^2},
  \]
  等号在
  \[
    u=\frac a{\sqrt{a^2+b^2}},\qquad
    v=\frac b{\sqrt{a^2+b^2}}
  \]
  时成立。因此所求长、宽为
  \[
    \boxed{
    \frac{2a^2}{\sqrt{a^2+b^2}},\qquad
    \frac{2b^2}{\sqrt{a^2+b^2}}}.
  \]
\end{enumerate}
\end{xsolution}

\paragraph{题 8}
\begin{xsolution}
设 $y_1,y_2$ 是两个满足同一边界条件的解，令 $z=y_1-y_2$，则
\[
  z''+p(x)z'+q(x)z=0,
  \qquad z(a)=z(b)=0.
\]
若 $z\not\equiv0$，则连续函数 $z$ 在 $(a,b)$ 内或者取得正的最大值，或者取得负的最小值。

若在 $x_0$ 处取得正最大值，则 $z'(x_0)=0$、$z''(x_0)\le0$，但微分方程给出
\[
  z''(x_0)=-q(x_0)z(x_0)>0,
\]
矛盾。若取得负最小值，则 $z'(x_0)=0$、$z''(x_0)\ge0$，而
\[
  z''(x_0)=-q(x_0)z(x_0)<0,
\]
也矛盾。故 $z\equiv0$，解唯一。
\end{xsolution}

\section*{8.4.2 练习题}

\paragraph{题 1}
\begin{xsolution}
对 $n$ 作归纳。$n=2$ 时就是下凸函数的定义。设结论对 $n-1$ 个点成立。令
\[
  \mu=\lambda_1+\cdots+\lambda_{n-1}=1-\lambda_n,
  \qquad
  y=\frac1\mu\sum_{i=1}^{n-1}\lambda_i x_i.
\]
由两点下凸性及归纳假设，
\begin{align*}
  f\left(\sum_{i=1}^n\lambda_i x_i\right)
  &=f(\mu y+\lambda_nx_n)\\
  &\le \mu f(y)+\lambda_nf(x_n)\\
  &\le \sum_{i=1}^{n-1}\lambda_i f(x_i)+\lambda_nf(x_n).
\end{align*}
故 Jensen 不等式对任意下凸函数成立，不需可微条件。
\end{xsolution}

\paragraph{题 2}
\begin{xsolution}
设 $z=\lambda x+(1-\lambda)y$，$0<\lambda<1$。因 $f$ 上凸且为正，
\[
  f(z)\ge \lambda f(x)+(1-\lambda)f(y)>0.
\]
取倒数并利用 $u\mapsto1/u$ 在 $(0,+\infty)$ 上下凸，得到
\begin{align*}
  \frac1{f(z)}
  &\le\frac1{\lambda f(x)+(1-\lambda)f(y)}\\
  &\le\frac\lambda{f(x)}+\frac{1-\lambda}{f(y)}.
\end{align*}
故 $1/f$ 下凸。

若把 $f$ 的条件改成下凸，则一般没有确定结论。例如 $f(x)=1+x^2$ 在 $\RR$ 上正且下凸，而
\[
  \left(\frac1{1+x^2}\right)''
  =\frac{2(3x^2-1)}{(1+x^2)^3}
\]
变号，所以 $1/f$ 既非全局下凸也非全局上凸。
\end{xsolution}

\paragraph{题 3}
\begin{xsolution}
令 $h=\max\{f,g\}$。对 $z=\lambda x+(1-\lambda)y$，有
\begin{align*}
  h(z)
  &=\max\{f(z),g(z)\}\\
  &\le\max\{\lambda f(x)+(1-\lambda)f(y),
             \lambda g(x)+(1-\lambda)g(y)\}\\
  &\le\lambda h(x)+(1-\lambda)h(y).
\end{align*}
故 $h$ 下凸。
\end{xsolution}

\paragraph{题 4}
\begin{xsolution}
设 $x<y$，$z=\lambda x+(1-\lambda)y$。记
\[
  a=f(x),\ b=f(y),\ c=g(x),\ d=g(y).
\]
由单调增加及非负性，$0\le a\le b$，$0\le c\le d$。由下凸性，
\[
  f(z)g(z)\le[\lambda a+(1-\lambda)b][\lambda c+(1-\lambda)d].
\]
而
\begin{align*}
 &\lambda ac+(1-\lambda)bd
 -[\lambda a+(1-\lambda)b][\lambda c+(1-\lambda)d]\\
 &=\lambda(1-\lambda)(b-a)(d-c)\ge0.
\end{align*}
故
\[
  f(z)g(z)\le\lambda f(x)g(x)+(1-\lambda)f(y)g(y),
\]
即 $fg$ 下凸。
\end{xsolution}

\paragraph{题 5}
\begin{xsolution}
设 $z=\lambda x+(1-\lambda)y$。由 $f$ 下凸且指数函数严格增加，
\[
  \ee^{f(z)}\le \ee^{\lambda f(x)+(1-\lambda)f(y)}.
\]
又因 $u\mapsto\ee^u$ 下凸，
\[
  \ee^{\lambda f(x)+(1-\lambda)f(y)}
  \le\lambda\ee^{f(x)}+(1-\lambda)\ee^{f(y)}.
\]
故 $F(x)=\ee^{f(x)}$ 下凸。
\end{xsolution}

\paragraph{题 6}
\begin{xsolution}
不一定。例如
\[
  f(u)=\ee^{-u},\qquad g(x)=x^2
\]
都是 $\RR$ 上的下凸函数，但
\[
  (f\circ g)(x)=\ee^{-x^2},
  \qquad
  (f\circ g)''(x)=(4x^2-2)\ee^{-x^2},
\]
二阶导数在 $x=0$ 附近为负，所以复合函数并非下凸函数。
\end{xsolution}

\paragraph{题 7}
\begin{xsolution}
由下凸函数的单侧导数性质，$f'_-$ 与 $f'_+$ 均单调增加。若所有单侧斜率均非负，则 $f$ 单调增加；若均非正，则 $f$ 单调减少。

余下情形中，负斜率与正斜率都出现。令
\[
  c=\sup\{x\in\RR\mid f'_+(x)<0\}.
\]
单侧导数的单调性表明：在 $x<c$ 的范围内各割线斜率非正，而在 $x>c$ 的范围内各割线斜率非负。由命题 8.4.1 的割线斜率刻画可得 $f$ 在 $(-\infty,c]$ 上单调减少，在 $[c,+\infty)$ 上单调增加。端点处若存在一段斜率为零，只需把 $c$ 取在该常值区间内即可。
\end{xsolution}

\paragraph{题 8}
\begin{xsolution}
由 $f''\ge0$，函数 $f$ 下凸，故任一点的切线都是支撑线。若某点 $x_0$ 有 $f'(x_0)>0$，则
\[
  f(x)\ge f(x_0)+f'(x_0)(x-x_0)\to+\infty
  \quad(x\to+\infty),
\]
与有界矛盾。若 $f'(x_0)<0$，令 $x\to-\infty$ 同样矛盾。因此 $f'(x)=0$ 对一切 $x$ 成立，故 $f$ 为常值函数。
\end{xsolution}

\paragraph{题 9}
\begin{xsolution}
因为 $f^{(n)}(x)>0$，由导数的 Darboux 性质及中值定理，$f^{(n-1)}$ 严格增加，而 $f^{(n-1)}(x_0)=0$，所以它在 $x_0$ 左负右正。逐级积分（或反复用中值定理）可知：每下降一次求导阶数，符号在左侧交替一次、右侧保持为正。

当 $n$ 为奇数时，从 $f^{(n)}$ 下降到 $f'$ 共 $n-1$ 次，$n-1$ 为偶数，因此 $f'(x)>0$ 对 $x\ne x_0$ 成立，而 $f'(x_0)=0$，故 $f$ 在 $(a,b)$ 上严格单调增加。

当 $n$ 为偶数时，从 $f^{(n)}$ 下降到 $f''$ 共 $n-2$ 次，$n-2$ 为偶数，因此 $f''(x)>0$ 对 $x\ne x_0$ 成立，而 $f''(x_0)=0$。故 $f'$ 严格单调增加，从而 $f$ 在 $(a,b)$ 上严格下凸。
\end{xsolution}

\paragraph{题 10}
\begin{xsolution}
任取闭区间 $[a,b]\subset(c,d)$。再取 $\alpha,\beta$ 使
\[
  c<\alpha<a<b<\beta<d.
\]
对任意 $a\le x_1<x_2\le b$，由下凸函数三点斜率不等式，
\[
  \frac{f(a)-f(\alpha)}{a-\alpha}
  \le\frac{f(x_2)-f(x_1)}{x_2-x_1}
  \le\frac{f(\beta)-f(b)}{\beta-b}.
\]
令两端常数分别为 $m,M$，则
\[
  \left|\frac{f(x_2)-f(x_1)}{x_2-x_1}\right|
  \le\max\{|m|,|M|\}.
\]
故 $f$ 在 $[a,b]$ 上满足 Lipschitz 条件。
\end{xsolution}

\paragraph{题 11}
\begin{xsolution}
先设 $f$ 下凸。固定 $c\in I$。由命题 8.4.3，有限单侧导数满足
\[
  f'_-(c)\le f'_+(c).
\]
任取
\[
  a\in[f'_-(c),f'_+(c)].
\]
若 $x<c$，割线斜率不大于 $f'_-(c)\le a$；若 $x>c$，割线斜率不小于 $f'_+(c)\ge a$。两种情形都给出
\[
  f(x)\ge f(c)+a(x-c).
\]

反之，设每个 $c$ 都存在这样的支撑线。对任意 $x_1<x_2<x_3$，在 $c=x_2$ 取支撑线斜率 $a$，则
\[
  \frac{f(x_2)-f(x_1)}{x_2-x_1}\le a
  \le\frac{f(x_3)-f(x_2)}{x_3-x_2}.
\]
由命题 8.4.1，$f$ 下凸。
\end{xsolution}

\paragraph{题 12}
\begin{xsolution}
对 $x>a$ 定义
\[
  q(x)=\frac{f(x)-f(a)}{x-a}.
\]
由凸函数三点斜率不等式，$q(x)$ 单调增加，故
\[
  L=\lim_{x\to+\infty}q(x)
\]
在扩充实数意义下存在（可能为 $+\infty$）。又
\[
  \frac{f(x)}x=\frac{x-a}{x}q(x)+\frac{f(a)}x,
\]
所以
\[
  \lim_{x\to+\infty}\frac{f(x)}x=L.
\]
这说明该极限一定有意义。
\end{xsolution}

\paragraph{题 13}
\begin{xsolution}
若 $f$ 不是常值函数，则存在 $x_1<x_2$ 使割线斜率
\[
  m=\frac{f(x_2)-f(x_1)}{x_2-x_1}\ne0.
\]
若 $m>0$，由下凸性，对 $x>x_2$ 有
\[
  \frac{f(x)-f(x_2)}{x-x_2}\ge m,
\]
故
\[
  \liminf_{x\to+\infty}\frac{f(x)}x\ge m>0,
\]
矛盾。若 $m<0$，向 $-\infty$ 作同样比较得到
\[
  \limsup_{x\to-\infty}\frac{f(x)}x\le m<0,
\]
也矛盾。因此所有割线斜率均为零，$f$ 为常值函数。
\end{xsolution}

\paragraph{题 14}
\begin{xsolution}
记公共值为 $K$。由 $f$ 在 $[a,b]$ 上凸，图像在端点弦下方，故
\[
  f(x)\le K,\qquad x\in[a,b].
\]
若 $x\in(a,c)$，可将 $c$ 写成 $x$ 与 $b$ 的严格凸组合：
\[
  c=\lambda x+(1-\lambda)b,\qquad0<\lambda<1.
\]
于是
\[
  K=f(c)\le\lambda f(x)+(1-\lambda)K,
\]
从而 $f(x)\ge K$。故 $f(x)=K$。对 $x\in(c,b)$ 用 $a$ 与 $x$ 表示 $c$ 同理。于是 $f\equiv K$。
\end{xsolution}

\paragraph{题 15}
\begin{xsolution}
用命题 8.4.1 的割线斜率刻画。因为 $f$ 在 $[a,c]$ 和 $[b,d]$ 上均下凸，所以各自区间内的割线斜率随位置右移而不减。特别是重叠区间 $[b,c]$ 把两套斜率次序连接起来。

任取 $x_1<x_2<x_3$ 于 $[a,d]$。若三点全在 $[a,c]$ 或全在 $[b,d]$，结论已有。若跨越两个区间，取任意 $u<v$ 于重叠区间 $[b,c]$，利用三点斜率不等式先在 $[a,c]$ 将左侧割线与 $[u,v]$ 的割线比较，再在 $[b,d]$ 将 $[u,v]$ 的割线与右侧割线比较，即得到
\[
  \frac{f(x_2)-f(x_1)}{x_2-x_1}
  \le
  \frac{f(x_3)-f(x_2)}{x_3-x_2}
\]
（端点落在重叠区间的情形直接取相应端点）。故任意三点都满足命题 8.4.1，$f$ 在 $[a,d]$ 上下凸。
\end{xsolution}

\paragraph{题 16}
\begin{xsolution}
设奇次多项式 $P$ 的次数为 $2m+1\ge3$。则 $P''$ 是次数 $2m-1$ 的奇次多项式，其首项系数非零，因此
\[
  P''(x)\to+\infty,\ -\infty
\]
在 $x\to+\infty$ 与 $x\to-\infty$ 时符号相反（次序由首项系数决定）。故 $P''$ 不可能在全实轴恒非负，也不可能恒非正。由二阶导数判据，$P$ 不可能在 $\RR$ 上为下凸或上凸函数。
\end{xsolution}

\paragraph{题 17}
\begin{xsolution}
因 $f''$ 无零点且导数具有 Darboux 性质，$f''$ 在 $(a,b)$ 上恒正或恒负。因此 $f'$ 严格单调。

对固定 $x_1<x_2$，Lagrange 中值定理中的 $\xi$ 必须满足
\[
  f'(\xi)=\frac{f(x_2)-f(x_1)}{x_2-x_1}.
\]
由于 $f'$ 为严格单调函数，同一个值至多有一个原像，所以 $\xi$ 唯一。
\end{xsolution}

\paragraph{题 18}
% CHECK-SOLUTION: 原书扫描页 250（物理页 267）明确写作“若曲线以该点为拐点且存在 f'''(x_0)，证明 f'''(x_0)=0”，但命题本身错误；反例 f(x)=x^3 在 x_0=0 为拐点而 f'''(0)=6。
\begin{xsolution}
按本书对拐点的定义，该命题不成立。例如
\[
  f(x)=x^3
\]
在 $x_0=0$ 两侧分别严格上凸、严格下凸（因为 $f''(x)=6x$ 变号），故 $(0,0)$ 是拐点；但
\[
  f'''(0)=6\ne0.
\]
因此题面所述命题为假。
\end{xsolution}

\paragraph{题 19}
\begin{xsolution}
不一定。例如 $f(x)=x^4$ 满足 $f'''(0)=0$，但
\[
  f''(x)=12x^2\ge0
\]
在原点两侧凸性不改变，所以 $(0,0)$ 不是拐点。
\end{xsolution}

\paragraph{题 20}
\begin{xsolution}
令 $g=f''$。由题设及 Peano 型 Taylor 公式，
\[
  f''(x)
  =\frac{f^{(n)}(x_0)}{(n-2)!}(x-x_0)^{n-2}
   +\littleo((x-x_0)^{n-2}).
\]
因 $f^{(n)}(x_0)\ne0$，在 $x_0$ 的充分小邻域内，$f''(x)$ 的符号由
\[
  f^{(n)}(x_0)(x-x_0)^{n-2}
\]
决定。曲线在两侧严格凸性不同，当且仅当该量在 $x_0$ 两侧变号，即当且仅当 $n-2$ 为奇数。故充分必要条件为
\[
  \boxed{n\text{ 为奇数}}.
\]
\end{xsolution}

\section*{8.5.3 练习题}

\paragraph{题 1}
\begin{xsolution}
令
\[
  F(x)=\ln x-\frac{2(x-1)}{x+1}.
\]
则 $F(1)=0$，且
\[
  F'(x)=\frac1x-\frac4{(x+1)^2}
  =\frac{(x-1)^2}{x(x+1)^2}>0\qquad(x>1).
\]
故 $F(x)>0$，即所证不等式成立。
\end{xsolution}

\paragraph{题 2}
\begin{xsolution}
函数 $h(x)=x\ln x$ 在 $(0,+\infty)$ 上严格下凸，因为 $h''(x)=1/x>0$。由 Jensen 不等式，且 $a\ne b$，
\[
  \frac{a\ln a+b\ln b}{2}>
  \frac{a+b}{2}\ln\frac{a+b}{2}.
\]
两边乘以 2 即得
\[
  a\ln a+b\ln b>(a+b)[\ln(a+b)-\ln2].
\]
\end{xsolution}

\paragraph{题 3}
\begin{xsolution}
令 $t=(x_2-x_1)/x_1>0$，则 $x_2/x_1=1+t$。例题 8.5.1 给出
\[
  \frac{t}{1+t}<\ln(1+t)<t.
\]
代回即为
\[
  \frac{x_2-x_1}{x_2}<\ln\frac{x_2}{x_1}
  <\frac{x_2-x_1}{x_1}.
\]

\textbf{另解}\quad 对 $\ln x$ 在区间 $[x_1,x_2]$ 使用 Lagrange 中值定理，存在 $\xi\in(x_1,x_2)$ 使
\[
  \frac{\ln x_2-\ln x_1}{x_2-x_1}=\frac1\xi.
\]
由于 $x_1<\xi<x_2$，有
\[
  \frac1{x_2}<\frac1\xi<\frac1{x_1},
\]
两边乘以 $x_2-x_1>0$ 即得所证双边不等式。
\end{xsolution}

\paragraph{题 4}
\begin{xsolution}
\begin{enumerate}[label=(\arabic*)]
  \item 函数 $u\mapsto a^u=\ee^{u\ln a}$ 为下凸函数（$a=1$ 时为常值），故 Jensen 不等式给出
  \[
    a^{(x_1+\cdots+x_n)/n}\le\frac1n\sum_{i=1}^n a^{x_i}.
  \]

  \item 在使 $x_i^p$ 为实数的通常非负变量范围内，$u\mapsto u^p$（$p>1$）下凸，因此
  \[
    \left(\frac{x_1+\cdots+x_n}{n}\right)^p
    \le\frac1n\sum_{i=1}^n x_i^p.
  \]

  \item 函数 $h(u)=u\ln u$ 在 $(0,+\infty)$ 上严格下凸。记
  \[
    S=\sum_{i=1}^n x_i,\qquad \bar x=\frac Sn.
  \]
  Jensen 不等式给出
  \[
    \frac1n\sum x_i\ln x_i\ge\bar x\ln\bar x.
  \]
  乘以 $n/S$ 后指数化，得到
  \[
    \frac Sn\le
    \exp\left(\frac1S\sum x_i\ln x_i\right)
    =\left(\prod x_i^{x_i}\right)^{1/S}.
  \]
\end{enumerate}
\end{xsolution}

\paragraph{题 5}
\begin{xsolution}
在 Bernoulli 不等式中令 $u=x-1>-1$。当 $0<\alpha<1$ 时
\[
  x^\alpha=(1+u)^\alpha\le1+\alpha u=\alpha x+1-\alpha,
\]
即
\[
  x^\alpha-\alpha x+\alpha-1\le0.
\]
当 $\alpha>1$ 时 Bernoulli 不等式反向，故结论也反向。
\end{xsolution}

\paragraph{题 6}
\begin{xsolution}
所求条件等价于
\[
  \alpha\le A_n\le\beta,
  \qquad
  A_n=\frac1{\ln(1+1/n)}-n.
\]
把 $n$ 暂看作连续变量 $x>0$。可证明 $A(x)$ 严格增加。事实上令 $t=1/x$，只需用
\[
  \ln(1+t)<\frac{t}{\sqrt{1+t}},\qquad t>0,
\]
而该式可由两边之差求导得到。于是
\[
  \inf_{n\ge1}A_n=A_1=\frac1{\ln2}-1.
\]
另一方面由
\[
  \ln(1+1/n)=\frac1n-\frac1{2n^2}+\bigO(n^{-3})
\]
得
\[
  \lim_{n\to\infty}A_n=\frac12.
\]
因此最大的 $\alpha$ 与最小的 $\beta$ 分别是
\[
  \boxed{\alpha=\frac1{\ln2}-1,\qquad \beta=\frac12}.
\]
\end{xsolution}

\paragraph{题 7}
\begin{xsolution}
依次记调和、几何、对数、算术、平方平均值为 $H,G,L,A,Q$。显然 $a<H$。由
\[
  (\sqrt a-\sqrt b)^2>0
\]
可得 $H<G$。

令 $u=\sqrt{b/a}>1$。不等式 $G<L$ 等价于
\[
  2\ln u<u-\frac1u.
\]
右减左的导数为 $(u-1)^2/u^2>0$，且在 $u=1$ 为 0，故成立。

由题 1 对 $x=b/a>1$ 应用
\[
  \ln x>\frac{2(x-1)}{x+1},
\]
得到 $L<A$。而
\[
  A<Q
\]
是严格的均方根－算术平均值不等式，$Q<b$ 由 $a^2<b^2$ 立即得到。因此
\[
  a<H<G<L<A<Q<b.
\]
\end{xsolution}

\paragraph{题 8}
\begin{xsolution}
先取有限的 $r<s$。若 $0<r<s$，令 $p=s/r>1$，对正数 $x_i^r$ 用加权幂平均（或 Jensen）得
\[
  \left(\sum\lambda_i x_i^r\right)^{1/r}
  \le
  \left(\sum\lambda_i x_i^s\right)^{1/s}.
\]
若 $r<s<0$，对 $1/x_i$ 使用前述正指数情形即可。若 $r<0<s$，分别与 $M_0=\prod x_i^{\lambda_i}$ 比较即可；$r=0$ 由极限连接。

当 $n>1$ 且 $x_i$ 不全相等时，所用 Jensen/幂平均不等式均为严格不等式，所以 $M_t$ 严格增加。最后
\[
  \lim_{t\to+\infty}M_t=\max_i x_i,
  \qquad
  \lim_{t\to-\infty}M_t=\min_i x_i,
\]
与题中的补充定义一致，故结论在整个扩充参数区间成立。
\end{xsolution}

\paragraph{题 9}
\begin{xsolution}
设 $0<p<1$，并记
\[
  A=\left(\sum x_k^p\right)^{1/p},
  \qquad
  B=\left(\sum y_k^p\right)^{1/p}.
\]
若 $A,B>0$，令 $u_k=x_k/A$，$v_k=y_k/B$，$\lambda=A/(A+B)$，$\mu=B/(A+B)$。由于 $t^p$ 在非负半轴上为上凸函数，
\[
  \sum(\lambda u_k+\mu v_k)^p
  \ge\lambda\sum u_k^p+\mu\sum v_k^p=1.
\]
因此
\[
  \left(\sum(x_k+y_k)^p\right)^{1/p}\ge A+B,
\]
即 Minkowski 不等式反向。$A=0$ 或 $B=0$ 时结论直接成立。
\end{xsolution}

\paragraph{题 10}
\begin{xsolution}
对非负数 $x,y$ 应用权为 $1/p,1/q$ 的加权算术－几何平均值不等式：
\[
  \frac xp+\frac yq
  \ge x^{1/p}y^{1/q}.
\]
这就是 Young 不等式。等号当且仅当 $x=y$。
\end{xsolution}

\paragraph{题 11}
\begin{xsolution}
\begin{enumerate}[label=(\arabic*)]
  \item 对两个数，Young 不等式正是权重为 $1/p$、$1/q$ 的加权 AM--GM。对多个有理权重反复使用二元形式即可得到广义 AM--GM；一般实权重再由连续性逼近得到。

  \item 记
  \[
    X=\left(\sum x_k^p\right)^{1/p},
    \qquad
    Y=\left(\sum y_k^q\right)^{1/q}.
  \]
  若 $XY=0$ 则显然。否则对每个 $k$ 将 Young 不等式用于
  \[
    \left(\frac{x_k}{X}\right)^p,
    \qquad
    \left(\frac{y_k}{Y}\right)^q,
  \]
  得
  \[
    \frac{x_ky_k}{XY}
    \le\frac1p\frac{x_k^p}{X^p}
       +\frac1q\frac{y_k^q}{Y^q}.
  \]
  对 $k$ 求和，右端为 $1/p+1/q=1$，故
  \[
    \sum x_ky_k\le XY,
  \]
  即 Hölder 不等式。
\end{enumerate}
\end{xsolution}

\paragraph{题 12}
\begin{xsolution}
由 $|f'|\le\varphi'$ 可知 $\varphi'\ge0$，且
\[
  (\varphi-f)'=\varphi'-f'\ge0,
  \qquad
  (\varphi+f)'=\varphi'+f'\ge0.
\]
因此对 $x\ge a$，利用 $f(a)=\varphi(a)$，
\[
  \varphi(x)-\varphi(a)\ge f(x)-f(a),
\]
同时
\[
  \varphi(x)-\varphi(a)\ge-[f(x)-f(a)].
\]
合并即得
\[
  |f(x)-f(a)|\le\varphi(x)-\varphi(a).
\]
\end{xsolution}

\paragraph{题 13}
\begin{xsolution}
因 $f'(x)>0$，$f$ 在 $[1,+\infty)$ 上严格增加，故 $f(x)\ge1$。于是
\[
  0<f'(x)=\frac1{x^2+f(x)^2}<\frac1{x^2+1}\qquad(x>1).
\]
令
\[
  \varphi(x)=1+\arctan x-\frac\pi4.
\]
则 $\varphi(1)=f(1)=1$ 且 $f'(x)<\varphi'(x)$，故由导数比较
\[
  1<f(x)<1+\arctan x-\frac\pi4<1+\frac\pi4.
\]
所以 $f$ 单调增加且有上界，极限
\[
  L=\lim_{x\to+\infty}f(x)
\]
存在。又由微分方程
\[
  L-1
  =\int_1^{+\infty}\frac{\dd t}{t^2+f(t)^2}.
\]
对 $t>1$ 有 $f(t)>1$，故被积函数严格小于 $1/(t^2+1)$。例如先在 $[1,2]$ 上积分即可得到严格差额，而在 $[2,+\infty)$ 上仍保持不大于，因而
\[
  L-1
  <\int_1^{+\infty}\frac{\dd t}{t^2+1}
  =\frac\pi4.
\]
所以
\[
  \boxed{\displaystyle\lim_{x\to+\infty}f(x)<1+\frac\pi4}.
\]
\end{xsolution}

\paragraph{题 14}
\begin{xsolution}
对每个 $i$，有
\[
  y_i=\sum_{j=1}^n a_{ij}x_j,
  \qquad a_{ij}\ge0,\quad \sum_j a_{ij}=1.
\]
由加权 AM--GM，
\[
  y_i\ge\prod_{j=1}^n x_j^{a_{ij}}.
\]
对 $i$ 相乘，并利用每列和也为 1，得到
\[
  \prod_{i=1}^n y_i
  \ge\prod_{j=1}^n x_j^{\sum_i a_{ij}}
  =\prod_{j=1}^n x_j.
\]
若某些 $x_j=0$，右端为 0，结论仍直接成立。
\end{xsolution}

\paragraph{题 15}
\begin{xsolution}
记
\[
  R_n(x)=\sin x-T_n(x).
\]
有 $R_n(0)=R_n'(0)=0$，并且
\[
  R_n''(x)=-R_{n-1}(x).
\]
从 $R_0(x)=\sin x-x<0$（$x>0$）开始归纳。若 $R_{n-1}$ 在 $x>0$ 上符号为 $(-1)^n$，则 $R_n''$ 的符号为 $(-1)^{n+1}$。由 $R_n'(0)=R_n(0)=0$ 两次利用单调性可知 $R_n$ 本身也具有符号 $(-1)^{n+1}$。因此 $n$ 为奇数时 $R_n>0$，$n$ 为偶数时 $R_n<0$，恰为题中两式。
\end{xsolution}

\paragraph{题 16}
\begin{xsolution}
要证
\[
  \ee^{-x}-\left(1-\frac xa\right)^a
  \le\frac{x^2}{a}\ee^{-x},
  \qquad 0\le x\le a,\ a\ge1.
\]
乘以 $\ee^x$ 后等价于
\[
  1-\ee^x\left(1-\frac xa\right)^a\le\frac{x^2}{a}.
\]
若 $x\ge\sqrt a$，右端不小于 1，而左端不大于 1，结论成立。

以下设 $0\le x<\sqrt a$。此时两边移项后都为正，只需证明
\[
  \ee^x\left(1-\frac xa\right)^a\ge1-\frac{x^2}{a}.
\]
取对数，令
\[
  H(x)=x+a\ln\left(1-\frac xa\right)
       -\ln\left(1-\frac{x^2}{a}\right).
\]
有 $H(0)=0$，直接化简得
\[
  H'(x)=
  \frac{x[(x-1)^2+a-1]}{(a-x)(a-x^2)}\ge0.
\]
故 $H(x)\ge0$，所需不等式成立。
\end{xsolution}

\section*{8.6.2 练习题}

\paragraph{题 1}
\begin{xsolution}
本题没有唯一答案；关键是使三幅图互相满足“斜率”和“斜率变化率”的关系。给出一个可核验的代表：取
\[
  f(x)=x^3-3x.
\]
则
\[
  f'(x)=3x^2-3,
  \qquad
  f''(x)=6x.
\]
因此 $f$ 是一条三次曲线，在 $x=\pm1$ 处有极值，在 $x=0$ 处有拐点；$f'$ 是顶点为 $(0,-3)$、零点为 $\pm1$ 的开口向上抛物线；$f''$ 是过原点、斜率为 6 的直线。这一组三图完整体现了原题要求的对应关系。
\end{xsolution}

\paragraph{题 2}
\begin{xsolution}
令
\[
  f(x)=\frac{x+1}{x^2+1}.
\]
计算得
\[
  f'(x)=-\frac{x^2+2x-1}{(x^2+1)^2},
\]
\[
  f''(x)=\frac{2(x-1)(x^2+4x+1)}{(x^2+1)^3}.
\]
所以三个拐点横坐标为
\[
  x=1,\qquad x=-2\pm\sqrt3.
\]
相应纵坐标分别为
\[
  1,\qquad \frac{1\pm\sqrt3}{4}.
\]
容易核对三点均满足
\[
  \boxed{y=\frac x4+\frac34},
\]
故三个拐点共线。

作图所需的其余信息为：$x=-1\pm\sqrt2$ 是两个驻点；$f(-1)=0$；$y=0$ 是水平渐近线；由 $f''$ 的三个简单零点可依次确定各凸性区间。上述信息与单调区间一起唯一确定其标准草图。
\end{xsolution}

\paragraph{题 3}
\begin{xsolution}
逐小题给出足以确定草图的定义域、极限、单调性、极值及渐近性质。
\begin{enumerate}[label=(\arabic*)]
  \item $y=(1+1/x)^x$，$x>0$。由 8.2.2 题 2 知它严格增加，且
  \[
    \lim_{x\to0+}y=1,
    \qquad
    \lim_{x\to+\infty}y=\ee.
  \]
  故值域为 $(1,\ee)$，$y=\ee$ 为水平渐近线。

  \item
  \[
    y=\frac1{1+x^2}\ee^{1/(1-x^2)},
    \qquad x\ne\pm1.
  \]
  它为偶函数。对 $x\ne0,\pm1$，
  \[
    \frac{y'}y=
    \frac{2x^3(3-x^2)}{(1-x^2)^2(1+x^2)}.
  \]
  因而在 $(0,1)$ 严格增加，从 $y(0)=\ee$ 增到 $+\infty$；在 $(1,\sqrt3)$ 从 $0$ 增到极大值
  \[
    y(\sqrt3)=\frac14\ee^{-1/2},
  \]
  再在 $(\sqrt3,+\infty)$ 严格减少到 0。负半轴由偶对称得到。$x=\pm1$ 为单侧性态不同的间断点，$y=0$ 为无穷远水平渐近线。

  \item
  \[
    y=\frac{x^4}{(1+x)^3},\qquad x\ne-1.
  \]
  有
  \[
    y'=\frac{x^3(x+4)}{(x+1)^4},
    \qquad
    y''=\frac{12x^2}{(x+1)^5}.
  \]
  故 $x=-4$ 为局部极大值点，值为 $-256/27$；$x=0$ 为局部极小值点，值为 0。$x=-1$ 是垂直渐近线。长除法给出
  \[
    y=x-3+\frac6{x+1}-\frac4{(x+1)^2}+\frac1{(x+1)^3},
  \]
  故 $y=x-3$ 是斜渐近线。$(-\infty,-1)$ 上为上凸，$(-1,+\infty)$ 上为下凸（$x=0$ 处二阶导数为零但凸性不变）。

  \item
  \[
    y=(1+x^2)\ee^{-x^2}
  \]
  为偶函数，
  \[
    y'=-2x^3\ee^{-x^2}.
  \]
  因而在 $(-\infty,0)$ 增加，在 $(0,+\infty)$ 减少，于 $x=0$ 取全局最大值 1，且两端趋于 0。又
  \[
    y''=2x^2(2x^2-3)\ee^{-x^2},
  \]
  拐点为 $x=\pm\sqrt{3/2}$。

  \item
  \[
    y=\frac{\ln x}{x},\qquad x>0.
  \]
  有
  \[
    y'=\frac{1-\ln x}{x^2},
    \qquad
    y''=\frac{2\ln x-3}{x^3}.
  \]
  因而从 $x\to0+$ 时的 $-\infty$ 增加，在 $x=\ee$ 取最大值 $1/\ee$，再减少到 $0+$；零点为 $x=1$，拐点为 $x=\ee^{3/2}$。

  \item
  \[
    y=x^{2/3}\ee^{-x},\qquad x\in\RR.
  \]
  在 $x\ne0$ 时
  \[
    y'=\ee^{-x}x^{-1/3}\left(\frac23-x\right).
  \]
  因而在 $(-\infty,0)$ 严格减少并由 $+\infty$ 降到 0；$x=0$ 为尖点式全局极小值点；在 $(0,2/3)$ 增加，在 $(2/3,+\infty)$ 减少，并在
  \[
    x=\frac23
  \]
  取局部极大值 $(2/3)^{2/3}\ee^{-2/3}$，最后趋于 0。

  \item
  \[
    y=\frac{(x-1)^3}{(x+1)^2},\qquad x\ne-1.
  \]
  有
  \[
    y'=\frac{(x-1)^2(x+5)}{(x+1)^3},
    \qquad
    y''=\frac{24(x-1)}{(x+1)^4}.
  \]
  所以 $x=-5$ 为局部极大值点，值为 $-27/2$；$x=1$ 为水平拐点，值为 0。$x=-1$ 是垂直渐近线。又
  \[
    y=x-5+\frac{12}{x+1}-\frac8{(x+1)^2},
  \]
  故 $y=x-5$ 为斜渐近线。

  \item
  \[
    y=\sqrt{\frac{x^3}{x-a}},\qquad a>0.
  \]
  定义域为
  \[
    (-\infty,0]\cup(a,+\infty).
  \]
  对 $y>0$，
  \[
    \frac{y'}y=\frac{2x-3a}{2x(x-a)}.
  \]
  因而左支在 $(-\infty,0]$ 严格减少到 $(0,0)$；右支在 $(a,3a/2)$ 减少，在 $(3a/2,+\infty)$ 增加，并在
  \[
    x=\frac{3a}{2},\qquad y=\frac{3\sqrt3}{2}a
  \]
  取最小值。$x=a$ 为右侧垂直渐近线；两端斜渐近线分别为
  \[
    y=-x-\frac a2\quad(x\to-\infty),
    \qquad
    y=x+\frac a2\quad(x\to+\infty).
  \]

  \item $y=x^x$，$x>0$。有
  \[
    y'=x^x(\ln x+1),
  \]
  故在 $(0,\ee^{-1})$ 减少，在 $(\ee^{-1},+\infty)$ 增加，于 $x=\ee^{-1}$ 取最小值 $\ee^{-1/\ee}$。又
  \[
    y''=x^x\left[(\ln x+1)^2+\frac1x\right]>0,
  \]
  故全程严格下凸；$x\to0+$ 时 $y\to1$，$x\to+\infty$ 时 $y\to+\infty$。

  \item $y=x^{1/x}$，$x>0$。有
  \[
    \frac{y'}y=\frac{1-\ln x}{x^2}.
  \]
  所以在 $(0,\ee)$ 增加，在 $(\ee,+\infty)$ 减少，于 $x=\ee$ 取最大值 $\ee^{1/\ee}$；且
  \[
    \lim_{x\to0+}y=0,
    \qquad
    \lim_{x\to+\infty}y=1.
  \]
\end{enumerate}
\end{xsolution}

\paragraph{题 4}
\begin{xsolution}
在第一象限中 $x,y>0$。若 $x=y$，显然得到直线
\[
  y=x.
\]
若 $x\ne y$，由 $x^y=y^x$ 得
\[
  \frac{\ln x}{x}=\frac{\ln y}{y}.
\]
令 $t=y/x>0$，$t\ne1$。则
\[
  x^{tx}=(tx)^x
  \quad\Longrightarrow\quad
  x^{t-1}=t,
\]
所以非对角支可参数化为
\[
  \boxed{x=t^{1/(t-1)},\qquad y=t^{t/(t-1)},\qquad t>0,\ t\ne1.}
\]
当 $t\to1$ 时 $(x,y)\to(\ee,\ee)$；$t\to0+$ 时趋向 $(+\infty,1)$；$t\to+\infty$ 时趋向 $(1,+\infty)$。交换 $t$ 与 $1/t$ 互换 $x,y$，故曲线关于 $y=x$ 对称。于是第一象限图像恰由直线 $y=x$ 与上述一条对称曲线组成。
\end{xsolution}

\paragraph{题 5}
\begin{xsolution}
\begin{enumerate}[label=(\arabic*)]
  \item $x=t^2-2$，$y=t(t^2-2)$。消去 $t$ 得
  \[
    y^2=x^2(x+2),\qquad x\ge-2.
  \]
  曲线关于 $x$ 轴对称；$(-2,0)$ 处为竖直切线；原点是二重点，两支切线斜率为 $\pm\sqrt2$。这些性质即可作图。

  \item $x=2t-t^2$，$y=3t-t^3$。令 $u=t-1$，则
  \[
    x=1-u^2,
    \qquad
    y=3x-1-u^3.
  \]
  因而
  \[
    y=3x-1\mp(1-x)^{3/2},\qquad x\le1.
  \]
  两支在 $(1,2)$ 相接并有共同切线斜率 3；在仿射坐标 $X=x-1$、$Y=y-2-3X$ 下满足 $Y^2=(-X)^3$，故为半三次尖点型曲线。

  \item
  \[
    x=\frac{3t}{1+t^3},\qquad y=\frac{3t^2}{1+t^3}.
  \]
  消去 $t$ 得著名的 Descartes 叶形线
  \[
    x^3+y^3=3xy.
  \]
  第一象限有一闭叶，在原点两支切线为两坐标轴；另一支具有斜渐近线
  \[
    x+y+1=0.
  \]

  \item $x=t-t^3$，$y=1-t^4$。由 $t\mapsto-t$ 可知曲线关于 $y$ 轴对称。它经过 $(0,1)$ 和原点；原点对应 $t=\pm1$，两支切线斜率分别为 $\pm2$。在
  \[
    t=\pm\frac1{\sqrt3}
  \]
  有竖直切线，对应点
  \[
    \left(\pm\frac{2}{3\sqrt3},\frac89\right).
  \]
  当 $|t|\to\infty$ 时 $y\to-\infty$，上述信息确定其对称草图。

  \item
  \[
    x=a(t-\sin t),\qquad y=a(1-\cos t)
  \]
  是摆线。每隔 $2\pi$ 周期平移一个拱；尖点为
  \[
    (2k\pi a,0),
  \]
  每一拱的最高点为
  \[
    ((2k+1)\pi a,2a).
  \]
  且
  \[
    \frac{\dd y}{\dd x}=\frac{\sin t}{1-\cos t}=\cot\frac t2.
  \]

  \item
  \[
    x=a(\cos t+t\sin t),\qquad
    y=a(\sin t-t\cos t)
  \]
  是圆的渐开线。写成极坐标关系可见
  \[
    r=a\sqrt{1+t^2},
    \qquad
    \arg(x+\ii y)=t-\arctan t.
  \]
  在 $t=0$ 从 $(a,0)$ 出发形成尖端，随后半径严格增加并绕圆向外展开；切线方向与 $(\cos t,\sin t)$ 同向。这些量给出标准渐开线草图。
\end{enumerate}
\end{xsolution}

\section*{8.7.3 练习题}

\paragraph{题 1}
\begin{xsolution}
\begin{enumerate}[label=(\arabic*)]
  \item 对
  \[
    f(x)=x^3-2x^2-4x-7
  \]
  用 Newton 迭代
  \[
    x_{n+1}=x_n-\frac{x_n^3-2x_n^2-4x_n-7}{3x_n^2-4x_n-4}.
  \]
  取 $x_0=4$，依次得到
  \[
    4,\ 3.67857143,\ 3.63287255,\ 3.63198114,\ 3.63198081,\ldots
  \]
  因而精确到 $0.0001$ 的根为
  \[
    \boxed{3.6320}.
  \]

  \item 写成 $f(x)=\sin x+x-1=0$，Newton 迭代为
  \[
    x_{n+1}=x_n-\frac{\sin x_n+x_n-1}{\cos x_n+1}.
  \]
  取 $x_0=0.5$，得到
  \[
    0.5,\ 0.51095795,\ 0.51097343,\ldots
  \]
  故根的近似值为
  \[
    \boxed{0.5110}.
  \]
\end{enumerate}
\end{xsolution}

\paragraph{题 2}
\begin{xsolution}
以下例子说明命题中的条件确实不能任意删去。

若导数在初值处为零，例如用 Newton 法解
\[
  f(x)=x^3-1=0,
\]
却取 $x_0=0$，则第一步就因 $f'(0)=0$ 而无法定义。

即使每一步都能定义，也可能不收敛。例如
\[
  f(x)=x^3-2x+2,\qquad x_0=0.
\]
Newton 映射为
\[
  N(x)=x-\frac{x^3-2x+2}{3x^2-2}.
\]
直接算得
\[
  N(0)=1,\qquad N(1)=0,
\]
故迭代在 $0,1$ 之间形成二周期而失败。这说明没有合适的单调性、凸性与初值控制时，即使方程有实根，Newton 法也可能不收敛到根。
\end{xsolution}

\paragraph{题 3}
\begin{xsolution}
二分法每一步把当前含根区间的长度减半。若以区间半长作为误差上界 $\varepsilon_n$，则
\[
  \varepsilon_{n+1}=\frac12\varepsilon_n.
\]
因此
\[
  \lim_{n\to\infty}\frac{\varepsilon_{n+1}}{\varepsilon_n}=\frac12\ne0,
\]
按本节定义二分法是一阶算法。
\end{xsolution}

\paragraph{题 4}
\begin{xsolution}
题 9 中求 $\sqrt A$ 的迭代式为
\[
  x_{n+1}=\frac{x_n(x_n^2+3A)}{3x_n^2+A},
  \qquad A>0,\quad x_0>0.
\]
令 $r=\sqrt A$，$t_n=x_n/r$，则
\[
  t_{n+1}=T(t_n),
  \qquad
  T(t)=\frac{t(t^2+3)}{3t^2+1}.
\]
直接计算得
\[
  T(t)-1=\frac{(t-1)^3}{3t^2+1},
  \qquad
  T(t)-t=\frac{2t(1-t^2)}{3t^2+1}.
\]
故当 $t>1$ 时 $1<T(t)<t$，当 $0<t<1$ 时 $t<T(t)<1$。因此 $\{t_n\}$ 单调有界并收敛，其极限只能为 $1$，即 $x_n\to r$。

更重要的是有精确误差关系
\[
  x_{n+1}-r
  =\frac{(x_n-r)^3}{3x_n^2+A}.
\]
于是
\[
  \lim_{n\to\infty}
  \frac{x_{n+1}-r}{(x_n-r)^3}
  =\frac1{3r^2+A}
  =\frac1{4A}\ne0.
\]
所以该算法为求平方根的三阶收敛算法。
\end{xsolution}

\paragraph{题 5}
\begin{xsolution}
令 $r=\sqrt[3]A>0$，并设 $t_n=x_n/r$。递推式化为
\[
  t_{n+1}=T(t_n),
  \qquad
  T(t)=\frac{t(t^3+2)}{2t^3+1}.
\]
直接因式分解：
\[
  T(t)-1=\frac{(t-1)^3(t+1)}{2t^3+1},
  \qquad
  T(t)-t=\frac{t(1-t^3)}{2t^3+1}.
\]
因此若 $t>1$，则 $1<T(t)<t$；若 $0<t<1$，则 $t<T(t)<1$。故从任意 $t_0>0$ 出发，数列单调有界并收敛，极限只能是 $1$。所以
\[
  \boxed{x_n\to r=\sqrt[3]A}.
\]

再令 $e_n=x_n-r$。由上面的精确因式分解，
\[
  e_{n+1}
  =r\frac{(e_n/r)^3(2+e_n/r)}{2(1+e_n/r)^3+1},
\]
故
\[
  \lim_{n\to\infty}\frac{e_{n+1}}{e_n^3}
  =\frac{2}{3r^2}\ne0.
\]
所以该算法为三阶收敛。
\end{xsolution}

\paragraph{题 6}
\begin{xsolution}
取
\[
  f(x)=x^k-A,
\]
Newton 公式给出
\[
  \boxed{x_{n+1}
  =x_n-\frac{x_n^k-A}{kx_n^{k-1}}
  =\frac1k\left((k-1)x_n+\frac{A}{x_n^{k-1}}\right)}.
\]
这里按通常实数 $k$ 次根问题取 $k>1$；$k=1$ 时一步即得精确值。令 $r=A^{1/k}$，$e_n=x_n-r$。Newton 法的标准 Taylor 误差式给出
\[
  e_{n+1}
  =\frac{f''(r)}{2f'(r)}e_n^2+\littleo(e_n^2)
  =\frac{k-1}{2r}e_n^2+\littleo(e_n^2).
\]
故当初值取在 $r$ 的充分小正邻域内时迭代收敛，且收敛阶为二阶。
\end{xsolution}

\paragraph{题 7}
\begin{xsolution}
要求 $1/A$，对方程
\[
  f(x)=\frac1x-A=0
\]
应用 Newton 法：
\[
  x_{n+1}
  =x_n-\frac{1/x_n-A}{-1/x_n^2}
  =\boxed{x_n(2-Ax_n)}.
\]
这个公式只含加法与乘法。

令
\[
  \delta_n=1-Ax_n.
\]
则有精确递推关系
\[
  \delta_{n+1}=\delta_n^2.
\]
若初值满足 $0<Ax_0<2$，即 $|\delta_0|<1$，则 $\delta_n\to0$，故 $x_n\to1/A$。又令 $e_n=x_n-1/A$，则
\[
  e_{n+1}=-A e_n^2,
\]
因此收敛恰为二阶。
\end{xsolution}

\section*{8.8.2 第一组参考题}

\paragraph{第一组参考题 1}
\begin{xsolution}
\begin{enumerate}[label=(\arabic*)]
  \item 记 $h(x)=f(x)+f'(x)$。给定 $\varepsilon>0$，存在 $X$，当 $x\ge X$ 时 $|h(x)|<\varepsilon$。于是
  \[
    [\ee^x(f(x)-\varepsilon)]'
    =\ee^x(h(x)-\varepsilon)<0,
  \]
  \[
    [\ee^x(f(x)+\varepsilon)]'
    =\ee^x(h(x)+\varepsilon)>0.
  \]
  对 $x\ge X$ 积分意义等价的单调性比较给出
  \[
    -\varepsilon+C_1\ee^{-x}\le f(x)
    \le\varepsilon+C_2\ee^{-x},
  \]
  其中常数由 $X$ 处的值决定。令 $x\to+\infty$，得
  \[
    -\varepsilon\le\liminf f(x)\le\limsup f(x)\le\varepsilon.
  \]
  再令 $\varepsilon\downarrow0$，即得 $f(x)\to0$。

  \item 设反面存在 $x_0$ 使 $f(x_0)>1$。由
  \[
    f+f'\le1
  \]
  得 $\ee^x(f-1)$ 单调不增。因此对 $x<x_0$，
  \[
    f(x)-1\ge\ee^{x_0-x}[f(x_0)-1]\to+\infty
    \quad(x\to-\infty),
  \]
  与 $f$ 有界矛盾。若存在 $f(x_0)<-1$，由 $f+f'\ge-1$ 知 $\ee^x(f+1)$ 单调不减，于 $x<x_0$ 得 $f(x)\to-\infty$，同样矛盾。故
  \[
    |f(x)|\le1\qquad(\forall x\in\RR).
  \]
\end{enumerate}
\end{xsolution}

\paragraph{第一组参考题 2}
\begin{xsolution}
由题设
\[
  \frac1x\ln\left(1+x+\frac{f(x)}x\right)\to3.
\]
有限两侧极限迫使括号趋于 1。若 $f(0)\ne0$，则 $f(x)/x$ 无界，不可能成立，故 $f(0)=0$。写 Taylor 展开
\[
  f(x)=f'(0)x+\frac12f''(0)x^2+\littleo(x^2).
\]
于是括号的极限为 $1+f'(0)$，故还必须有 $f'(0)=0$。此时
\[
  1+x+\frac{f(x)}x
  =1+\left(1+\frac12f''(0)\right)x+\littleo(x).
\]
因此
\[
  3=1+\frac12f''(0),
  \qquad f''(0)=4.
\]
最后
\[
  1+\frac{f(x)}x=1+2x+\littleo(x),
\]
故
\[
  \boxed{f(0)=0,\quad f'(0)=0,\quad f''(0)=4,\quad
  \lim_{x\to0}\left(1+\frac{f(x)}x\right)^{1/x}=\ee^2}.
\]
\end{xsolution}

\paragraph{第一组参考题 3}
\begin{xsolution}
由正数列 $x_n\to0$ 且
\[
  x_{n+1}=x_n+Ax_n^k+\littleo(x_n^k),
\]
必有 $A<0$；否则在充分大的 $n$ 上 $x_{n+1}>x_n$，与正数列趋于 0 矛盾。

令
\[
  y_n=x_n^{-(k-1)}\to+\infty.
\]
则
\begin{align*}
  y_{n+1}-y_n
  &=x_n^{-(k-1)}
  \left[
    \left(1+Ax_n^{k-1}+\littleo(x_n^{k-1})\right)^{-(k-1)}-1
  \right]\\
  &\to-(k-1)A>0.
\end{align*}
由 Stolz 定理，
\[
  \frac{y_n}{n}\to-(k-1)A.
\]
因此应取
\[
  \boxed{\alpha=k-1},
\]
且
\[
  \boxed{\lim_{n\to\infty}n x_n^{k-1}
  =-\frac1{(k-1)A}>0}.
\]
\end{xsolution}

\paragraph{第一组参考题 4}
\begin{xsolution}
严格单调函数显然没有内点极值。

反之，设 $f$ 在 $(a,b)$ 内没有极值点。先证 $f$ 为一一函数。若存在 $x<y$ 使 $f(x)=f(y)$，则连续函数在 $[x,y]$ 上取得最大、最小值。若 $f$ 在该区间恒定，则任一内点都是极值点；若不恒定，则最大值或最小值中至少有一个与端点公共值不同，从而必在 $(x,y)$ 内取得，仍产生极值点。均矛盾。因此 $f$ 一一。

连续的一一实函数在区间上必严格单调，故 $f$ 在 $[a,b]$ 上严格单调。
\end{xsolution}

\paragraph{第一组参考题 5}
\begin{xsolution}
只需证明对 $u,v\ge0$，
\[
  (u+v)^p\le u^p+v^p,
  \qquad0<p<1.
\]
若 $u=0$ 或 $v=0$ 显然。设 $u>0$，令 $t=v/u\ge0$，只需证
\[
  (1+t)^p\le1+t^p.
\]
令 $F(t)=1+t^p-(1+t)^p$，则 $F(0)=0$，且
\[
  F'(t)=p[t^{p-1}-(1+t)^{p-1}]>0
\]
因为 $p-1<0$。故 $F(t)\ge0$。取 $u=|a|$、$v=|b|$，再用三角不等式，
\[
  |a+b|^p\le(|a|+|b|)^p\le|a|^p+|b|^p.
\]
\end{xsolution}

\paragraph{第一组参考题 6}
\begin{xsolution}
令 $t=1/n>0$，并记
\[
  r=1-\frac{\ln(1+t)}t>0.
\]
则
\[
  \frac1\ee\left[\ee-(1+t)^{1/t}\right]=1-\ee^{-r}.
\]
先证上界。由 $1-\ee^{-r}<r$ 以及例题 8.5.1 的强化式
\[
  \ln(1+t)>\frac{2t}{2+t}
\]
得
\[
  r<\frac{t}{2+t},
\]
所以
\[
  1-\ee^{-r}<\frac{t}{2+t}=\frac1{2n+1}.
\]

再证下界。要证
\[
  1-\ee^{-r}>\frac{t}{2(1+t)},
\]
等价于
\[
  r>\ln\frac{1+t}{1+t/2}.
\]
令
\[
  H(t)=1+\ln(1+t/2)-\left(1+\frac1t\right)\ln(1+t).
\]
有 $H(0+)=0$，并且
\[
  H'(t)=\frac{(t+2)\ln(1+t)-2t}{t^2(t+2)}>0,
\]
最后一步仍由 $\ln(1+t)>2t/(2+t)$。故 $H(t)>0$，下界成立。恢复 $t=1/n$，得到
\[
  \boxed{
  \frac{\ee}{2n+2}<\ee-\left(1+\frac1n\right)^n
  <\frac{\ee}{2n+1}}.
\]
\end{xsolution}

\paragraph{第一组参考题 7}
\begin{xsolution}
令
\[
  F(x)=3\ln\frac{\sin x}{x}-\ln\cos x,
  \qquad0<x<\frac\pi2.
\]
有 $F(0+)=0$。化简导数：
\[
  F'(x)=
  \frac{x(1+2\cos^2x)-3\sin x\cos x}{x\sin x\cos x}.
\]
记分子为 $N(x)$。则 $N(0)=0$，且
\[
  N'(x)=4\sin x(\sin x-x\cos x).
\]
而 $G(x)=\sin x-x\cos x$ 满足 $G(0)=0$、$G'(x)=x\sin x>0$，故 $G(x)>0$，从而 $N'(x)>0$。所以 $F'(x)>0$、$F(x)>0$。指数化即得
\[
  \left(\frac{\sin x}{x}\right)^3>\cos x.
\]
\end{xsolution}

\paragraph{第一组参考题 8}
\begin{xsolution}
令
\[
  F(x)=2\sin x+\tan x-3x.
\]
则 $F(0)=0$，且
\begin{align*}
  F'(x)
  &=2\cos x+\operatorname{sec}^2x-3\\
  &=\frac{(1-\cos x)^2(2\cos x+1)}{\cos^2x}>0
\end{align*}
对 $0<x<\pi/2$ 成立。因此 $F(x)>0$，即所证不等式。
\end{xsolution}

\paragraph{第一组参考题 9}
\begin{xsolution}
先证下界。令
\[
  h(x)=\sin\pi x-\pi x(1-x).
\]
函数关于 $x=1/2$ 对称，只需看 $0<x<1/2$。此时
\[
  h'(x)=\pi[\cos\pi x-1+2x].
\]
因 $\cos t$ 在 $[0,\pi/2]$ 上严格上凸，图像严格高于连接端点的弦
\[
  \cos t>1-\frac{2t}{\pi},
\]
故 $h'(x)>0$，于是 $h(x)>h(0)=0$，即
\[
  \frac{\sin\pi x}{x(1-x)}>\pi.
\]

再证上界，令
\[
  H(x)=4x(1-x)-\sin\pi x.
\]
仍只需看 $[0,1/2]$。有 $H(0)=H(1/2)=0$，
\[
  H'=4-8x-\pi\cos\pi x,
  \quad H''=-8+\pi^2\sin\pi x,
  \quad H'''=\pi^3\cos\pi x>0.
\]
故 $H''$ 严格增加并只变号一次，$H'$ 先减后增；结合 $H'(0)=4-\pi>0$、$H'(1/2)=0$ 以及 $H(1/2)-H(0)=0$，可知 $H'$ 在内部只从正变负一次，故 $H$ 先增后减并保持非负。于是
\[
  \sin\pi x\le4x(1-x).
\]
等号在 $x=1/2$ 成立。综上即得题中双边不等式。
\end{xsolution}

\paragraph{第一组参考题 10}
\begin{xsolution}
记
\[
  F(x)=\sin x-\frac2\pi x-\frac1{12\pi}x(\pi^2-4x^2).
\]
直接验证 $F(0)=F(\pi/2)=0$，而
\[
  F''(x)=\frac{2x}{\pi}-\sin x\le0
\]
由 Jordan 不等式成立。因此 $F$ 在 $[0,\pi/2]$ 上为上凸函数。上凸函数的图像在连接两个端点的弦之上，而该弦就是 $y=0$，故 $F(x)\ge0$，即得所证强化不等式。
\end{xsolution}

\paragraph{第一组参考题 11}
\begin{xsolution}
记
\[
  S=x_1+\cdots+x_n,
  \quad y_i=\frac{x_i}{S},
  \quad z_i=\frac{1+x_i}{n+S}.
\]
则 $\sum y_i=\sum z_i=1$，且
\[
  z_i=\frac{S}{n+S}y_i+\frac{n}{n+S}\frac1n.
\]
由于 $\ln$ 为上凸函数，
\[
  \sum\ln z_i
  \ge\frac{S}{n+S}\sum\ln y_i
   +\frac{n}{n+S}\,n\ln\frac1n.
\]
而 AM--GM 给出
\[
  \sum\ln y_i\le n\ln\frac1n.
\]
因此右端不小于 $\sum\ln y_i$，从而
\[
  \prod z_i\ge\prod y_i.
\]
展开定义正是题中不等式。两处等号同时成立当且仅当
\[
  x_1=x_2=\cdots=x_n.
\]
\end{xsolution}

\paragraph{第一组参考题 12}
\begin{xsolution}
两边取对数。所需比较等价于
\[
  \frac{\ln\sin x}{\sin x}
  \quad\text{与}\quad
  \frac{\ln\cos x}{\cos x}
\]
的大小。令
\[
  \phi(t)=\frac{\ln t}{t},\qquad0<t\le1.
\]
则
\[
  \phi'(t)=\frac{1-\ln t}{t^2}>0,
\]
所以 $\phi$ 严格增加。当 $0<x<\pi/4$ 时 $\sin x<\cos x$，故
\[
  \frac{\ln\sin x}{\sin x}<\frac{\ln\cos x}{\cos x},
\]
等价于 $(\sin x)^{\cos x}<(\cos x)^{\sin x}$。当 $\pi/4<x<\pi/2$ 时次序反向，故不等式也反向。
\end{xsolution}

\paragraph{第一组参考题 13}
\begin{xsolution}
\begin{enumerate}[label=(\arabic*)]
  \item 对两个正数
  \[
    X=t^{1/p-1}a,
    \qquad Y=t^{1/p}b
  \]
  应用权分别为 $1/p$ 和 $1-1/p$ 的加权 AM--GM：
  \[
    \frac Xp+\left(1-\frac1p\right)Y
    \ge X^{1/p}Y^{1-1/p}.
  \]
  右端的 $t$ 次幂恰相消，故得到题中 (1)。

  \item 令
  \[
    F(t)=\frac{a^p}{t^{p-1}}+
          \frac{b^p}{(1-t)^{p-1}},\qquad0<t<1.
  \]
  有
  \[
    F'(t)=-(p-1)\frac{a^p}{t^p}
    +(p-1)\frac{b^p}{(1-t)^p}.
  \]
  唯一驻点为 $t=a/(a+b)$，且 $F$ 在其左侧减少、右侧增加，所以该点为全局最小值点。代入得到
  \[
    F_{\min}=(a+b)^p.
  \]
  这就是题中 (2)。
\end{enumerate}
\end{xsolution}

\paragraph{第一组参考题 14}
\begin{xsolution}
写成
\[
  p(x)=C(x-a)(x-b)(x-c),\qquad C\ne0,
\]
其中第三个根 $c$ 为实数。$p$ 在 $[a,b]$ 上不变号，当且仅当第三个根不在开区间 $(a,b)$ 中，即
\[
  (a-c)(b-c)\ge0.
\]
另一方面
\[
  p'(a)p'(b)
  =-C^2(a-b)^2(a-c)(b-c).
\]
故
\[
  p'(a)p'(b)\le0
  \quad\Longleftrightarrow\quad
  (a-c)(b-c)\ge0,
\]
恰与不变号条件等价。
\end{xsolution}

\paragraph{第一组参考题 15}
\begin{xsolution}
设右端多项式统一记为 $P$。若 $P=g\circ g$，则必有交换关系
\[
  g(P(x))=P(g(x)).
\]
因此 $g$ 把 $P$ 的不动点映到 $P$ 的不动点。
\begin{enumerate}[label=(\arabic*)]
  \item $P(x)=-x^3+x+1$ 的唯一不动点为 $1$。故 $g(1)=1$。链式法则给出
  \[
    P'(1)=[g'(1)]^2\ge0,
  \]
  但 $P'(1)=-2$，矛盾。

  \item $P(x)=-x^3+x^2+1$ 的唯一不动点仍为 $1$，同理应有 $P'(1)=[g'(1)]^2\ge0$，但 $P'(1)=-1$，矛盾。

  \item $P(x)=x^2-3x+3$ 的不动点为 $1,3$。若 $g(1)=1$，则 $P'(1)=[g'(1)]^2\ge0$，但 $P'(1)=-1$，矛盾。因此 $g(1)=3$，又由 $g(g(1))=P(1)=1$ 得 $g(3)=1$。于是
  \[
    P'(1)=g'(3)g'(1)=P'(3),
  \]
  但 $P'(1)=-1$、$P'(3)=3$，矛盾。
\end{enumerate}
故三式均不存在全实轴可微解。
\end{xsolution}

\paragraph{第一组参考题 16}
\begin{xsolution}
因端点值均为 0 而最小值为 $-1$，存在 $c\in(0,1)$ 使
\[
  f(c)=-1,\qquad f'(c)=0.
\]
对区间 $[0,c]$ 在 $c$ 点作二阶 Taylor 公式，存在 $\xi_1\in(0,c)$ 使
\[
  0=f(0)=f(c)+\frac12f''(\xi_1)c^2,
\]
故
\[
  f''(\xi_1)=\frac2{c^2}.
\]
同理存在 $\xi_2\in(c,1)$ 使
\[
  f''(\xi_2)=\frac2{(1-c)^2}.
\]
$c$ 与 $1-c$ 至少一个不超过 $1/2$，所以相应的二阶导数至少为 8。故存在所需 $\xi$。
\end{xsolution}

\paragraph{第一组参考题 17}
% CHECK-SOLUTION: 原书扫描页 275（物理页 292）明确写作“存在 ξ 使 f''(ξ)=3”，但该命题为假。反例：f(x)=x^2+x^3（x≤0），f(x)=x^2（x≥0）是 [-1,1] 上二阶可微函数，满足全部给定值，而 f'' 在左半区为 2+6x、右半区恒为 2，从不等于 3。
\begin{xsolution}
原题按所印条件不能成立。取
\[
  f(x)=
  \begin{cases}
    x^2+x^3,&-1\le x\le0,\\
    x^2,&0\le x\le1.
  \end{cases}
\]
则两侧在 $0$ 处的一、二阶导数均接合，故 $f$ 在 $[-1,1]$ 上二阶可微，并且
\[
  f(0)=f'(0)=0,\qquad f(1)=1,\qquad f(-1)=0.
\]
但
\[
  f''(x)=
  \begin{cases}
    2+6x,&x<0,\\
    2,&x\ge0,
  \end{cases}
\]
在 $(-1,1)$ 内从不取值 3。因此题面所述结论为假。
\end{xsolution}

\paragraph{第一组参考题 18}
\begin{xsolution}
反设 $g$ 在 $(a,b)$ 内无零点。先看端点：由 $f(a)=0$ 与 $W(f,g)(a)\ne0$，
\[
  W(f,g)(a)=-g(a)f'(a)\ne0,
\]
所以 $g(a)\ne0$；同理 $g(b)\ne0$。因此在反设下 $g$ 在整个 $[a,b]$ 上处处非零，可以定义
\[
  h(x)=\frac{f(x)}{g(x)}.
\]
此时 $h$ 在 $[a,b]$ 上连续、在 $(a,b)$ 上可微，并且
\[
  h(a)=h(b)=0.
\]
由 Rolle 定理，存在 $\xi\in(a,b)$ 使 $h'(\xi)=0$。但
\[
  h'(x)=\frac{f'g-fg'}{g^2}
  =-\frac{W(f,g)}{g^2}\ne0,
\]
矛盾。故 $g$ 必在 $(a,b)$ 中有零点。
\end{xsolution}

\section*{8.8.2 第二组参考题}

\paragraph{第二组参考题 1}
\begin{xsolution}
因 $f$ 单调减少且可微，有 $f'(x)\le0$；题设 $0<f(x)<|f'(x)|$ 进一步给出
\[
  -\frac{f'(x)}{f(x)}>1.
\]
当 $0<x<1$ 时，在 $[x,1/x]$ 上积分，得到
\[
  \ln\frac{f(x)}{f(1/x)}
  =\int_x^{1/x}-\frac{f'(t)}{f(t)}\,\dd t
  >\frac1x-x.
\]
又令
\[
  H(x)=\frac1x-x+2\ln x,
  \qquad 0<x\le1.
\]
则 $H(1)=0$，且
\[
  H'(x)=-\frac{(1-x)^2}{x^2}<0\qquad(0<x<1),
\]
故 $H(x)>0$，即
\[
  \frac1x-x>-2\ln x=\ln\frac1{x^2}.
\]
于是
\[
  \frac{f(x)}{f(1/x)}>\frac1{x^2},
\]
也就是
\[
  \boxed{x f(x)>\frac1x f\left(\frac1x\right)}.
\]
\end{xsolution}

\paragraph{第二组参考题 2}
\begin{xsolution}
对 $n$ 作归纳。$n=0$ 时结论就是 $g(0)=f'(0)$。设结论对 $n-1$ 已成立，则 $g^{(n-1)}$ 在 $0$ 连续。

对 $x\ne0$，由 $g(x)=f(x)x^{-1}$ 的 Leibniz 公式，
\[
  g^{(n)}(x)
  =\frac{n!}{x^{n+1}}
   \sum_{k=0}^n\frac{(-1)^{n-k}}{k!}x^k f^{(k)}(x).
\tag{1}
\]
由 $f^{(n+1)}(0)$ 存在，对 $0\le k\le n$ 有 Peano 展开式
\[
  f^{(k)}(x)
  =\sum_{j=0}^{n+1-k}
    \frac{f^{(k+j)}(0)}{j!}x^j
   +\littleo(x^{n+1-k}).
\]
代入 (1)。对 $1\le r\le n$，$x^r$ 的系数含组合和
\[
  \sum_{k=0}^r\frac{(-1)^{n-k}}{k!(r-k)!}=0,
\]
而常数项因 $f(0)=0$ 也为 0。对 $r=n+1$，由于求和只到 $k=n$，系数为
\[
  \sum_{k=0}^n\frac{(-1)^{n-k}}{k!(n+1-k)!}
  =\frac1{(n+1)!}.
\]
故
\[
  g^{(n)}(x)
  =\frac1{n+1}f^{(n+1)}(0)+\littleo(1)
  \qquad(x\to0,\ x\ne0).
\]
因此 $g^{(n)}(x)$ 在穿孔邻域中有有限极限
\[
  L=\frac1{n+1}f^{(n+1)}(0).
\]
由归纳假设 $g^{(n-1)}$ 在 $0$ 连续，再应用导数极限定理，$g^{(n-1)}$ 在 $0$ 可微且其导数等于 $L$。于是
\[
  \boxed{g^{(n)}(0)=\frac1{n+1}f^{(n+1)}(0)},
\]
并且 $g^{(n)}(x)\to g^{(n)}(0)$，即 $g^{(n)}$ 在 $0$ 连续。归纳完成。
\end{xsolution}

\paragraph{第二组参考题 3}
\begin{xsolution}
把
\[
  F(x)=f(x)^n
\]
看成 $n$ 个因子 $f(x)$ 的乘积。对 $0\le k\le n-1$，用 Leibniz 公式求 $F^{(k)}(a)$。每一项都是形如
\[
  C\,f^{(j_1)}(a)f^{(j_2)}(a)\cdots f^{(j_n)}(a),
  \qquad j_1+\cdots+j_n=k
\]
的乘积。因为 $k<n$，至少有一个 $j_i=0$，故该项含因子 $f(a)=0$。所以所有项均为 0，即
\[
  \boxed{F^{(k)}(a)=0,\qquad k=0,1,\ldots,n-1.}
\]

若完全去掉 $f^{(n)}(a)$ 存在这一正则性条件，结论可能连“导数存在”都失效。例如 $n\ge2$、$a=0$ 时取
\[
  f(x)=|x|^{1/n}.
\]
则 $f(0)=0$，但
\[
  F(x)=f(x)^n=|x|
\]
在 $0$ 处连一阶导数都不存在，故原条件不能无条件删去。
\end{xsolution}

\paragraph{第二组参考题 4}
\begin{xsolution}
记
\[
  P_n(x)=\sum_{k=0}^n\frac{x^k}{k!}.
\]
先注意 Taylor 公式给出：对任意 $x$，存在介于 $0$ 与 $x$ 之间的 $\xi$，使
\[
  \ee^x=P_n(x)+\frac{\ee^\xi x^{n+1}}{(n+1)!}.
\]
因此当 $x<0$ 时，
\[
  P_{2m}(x)>\ee^x>P_{2m+1}(x).
\tag{1}
\]

\begin{enumerate}[label=(\arabic*)]
  \item 若 $x\ge0$，$P_{2m}(x)>0$ 显然；若 $x<0$，由 (1) 有
  \[
    P_{2m}(x)>\ee^x>0.
  \]

  \item
  \[
    P_{2m+1}'(x)=P_{2m}(x)>0,
  \]
  故 $P_{2m+1}$ 在 $\RR$ 上严格增加。又其首项次数为奇数且首项系数为正，所以
  \[
    P_{2m+1}(-\infty)=-\infty,
    \qquad P_{2m+1}(+\infty)=+\infty.
  \]
  因而恰有一个实零点。

  \item 记 $P_{2n+1}$ 的唯一零点为 $x_n<0$，令 $z_n=-x_n>0$。Taylor 的积分余项写成
  \[
    \ee^{-z}
    =P_{2n+1}(-z)
     +\frac{z^{2n+2}}{(2n+1)!}
      \int_0^1(1-t)^{2n+1}\ee^{-tz}\,\dd t.
  \]
  令 $z=z_n$ 并乘以 $\ee^{z_n}$，得到
  \[
    1=\frac{z_n^{2n+2}}{(2n+1)!}
      \int_0^1(1-t)^{2n+1}\ee^{(1-t)z_n}\,\dd t
    >\frac{z_n^{2n+2}}{(2n+2)!}.
  \]
  因而 $z_n<2n+3$。于是
  \begin{align*}
    P_{2n+3}(x_n)
    &=\frac{x_n^{2n+2}}{(2n+2)!}
      \left(1+\frac{x_n}{2n+3}\right)>0.
  \end{align*}
  而 $P_{2n+3}$ 严格增加，故其零点满足
  \[
    x_{n+1}<x_n.
  \]
  所以 $\{x_n\}$ 严格减少。若它有有限下界，则 $x_n\to L\in\RR$；但 $P_m\to\ee^x$ 在每个有界闭区间上一致成立，故
  \[
    0=P_{2n+1}(x_n)\to\ee^L>0,
  \]
  矛盾。因此
  \[
    \boxed{x_n\to-\infty}.
  \]

  \item 这正是 (1)。

  \item 当 $x>0$ 时 Taylor 余项为正，故
  \[
    \ee^x>P_n(x).
  \]
  又由二项式定理
  \[
    \left(1+\frac xn\right)^n
    =\sum_{k=0}^n\binom nk\frac{x^k}{n^k}
    \le\sum_{k=0}^n\frac{x^k}{k!}=P_n(x),
  \]
  因为
  \[
    \frac1{n^k}\binom nk
    =\frac1{k!}\prod_{j=0}^{k-1}\left(1-\frac jn\right)
    \le\frac1{k!}.
  \]

  \item Taylor 余项满足
  \[
    |\ee^x-P_n(x)|
    \le \ee^{|x|}\frac{|x|^{n+1}}{(n+1)!}\to0,
  \]
  因而对每个固定 $x\in\RR$，
  \[
    \boxed{P_n(x)\to\ee^x}.
  \]
\end{enumerate}
\end{xsolution}

\paragraph{第二组参考题 5}
\begin{xsolution}
令
\[
  d_n=x_n-x_{n-1}\qquad(n\ge1).
\]
由递推式
\[
  x_{n+1}=\frac{(2n-1)x_n+x_{n-1}}{2n}
\]
得到
\[
  d_{n+1}=x_{n+1}-x_n=-\frac1{2n}d_n.
\]
因为 $d_1=b-a$，故
\[
  d_n=\frac{(-1)^{n-1}}{2^{n-1}(n-1)!}(b-a).
\]
于是
\begin{align*}
  \lim_{n\to\infty}x_n
  &=a+\sum_{n=1}^\infty d_n\\
  &=a+(b-a)\sum_{j=0}^\infty\frac{(-1)^j}{2^j j!}\\
  &=\boxed{a+(b-a)\ee^{-1/2}}.
\end{align*}
\end{xsolution}

\paragraph{第二组参考题 6}
\begin{xsolution}
对任意初值 $z_0\in(0,\pi/2]$，递推 $z_{n+1}=\sin z_n$ 产生严格减少并趋于 0 的正数列。例题 8.1.10 已证明其共同的渐近式
\[
  z_n\sim\sqrt{\frac3n}.
\]
分别应用于 $\{x_n\}$ 与 $\{y_n\}$，有
\[
  x_n\sim\sqrt{\frac3n},
  \qquad
  y_n\sim\sqrt{\frac3n}.
\]
因此
\[
  \boxed{\lim_{n\to\infty}\frac{x_n}{y_n}=1}.
\]
\end{xsolution}

\paragraph{第二组参考题 7}
\begin{xsolution}
设
\[
  y=\frac{p(x)}{q(x)},
  \quad
  p(x)=ax^2+2bx+c,
  \quad
  q(x)=\alpha x^2+2\beta x+\gamma.
\]
若 $\alpha=0$，直接约去多项式部分后，剩余部分至多是常数除以一次式，其二阶导数至多有一个零点；故存在三个拐点时必有 $\alpha\ne0$。作横坐标平移 $t=x+\beta/\alpha$，并从 $y$ 中减去常数 $a/\alpha$，可把函数化成
\[
  Y(t)=\frac{mt+n}{t^2+d}
\]
的形式。若 $d=0$，直接求导可见 $Y''$ 的分子至多一次，仍不可能有三个拐点；因此在题设情形 $d\ne0$。

计算
\[
  Y''(t)=
  \frac{2\{m(t^3-3dt)+n(3t^2-d)\}}{(t^2+d)^3}.
\]
记分子花括号内的三次多项式为
\[
  S(t)=m(t^3-3dt)+n(3t^2-d).
\]
有恒等式
\[
  mt+n
  =\frac{mt+3n}{4d}(t^2+d)-\frac1{4d}S(t).
\]
因此在每一个拐点参数 $t_i$（即 $S(t_i)=0$）处，
\[
  Y(t_i)=\frac{mt_i+3n}{4d}.
\]
这说明三个拐点都落在同一直线
\[
  Y=\frac{m}{4d}t+\frac{3n}{4d}
\]
上。平移坐标不改变共线性，故原函数的三个拐点必共线。
\end{xsolution}

\paragraph{第二组参考题 8}
\begin{xsolution}
§8.4 的下凸定义显然推出 Jensen 中点不等式，只需证明反向。

设 $f$ 连续且满足
\[
  f\left(\frac{x+y}{2}\right)\le\frac{f(x)+f(y)}2.
\]
反复使用中点不等式，可归纳得到：对每个二进有理数
\[
  \lambda=\frac{k}{2^m}\in[0,1],
\]
都有
\[
  f(\lambda x+(1-\lambda)y)
  \le\lambda f(x)+(1-\lambda)f(y).
\]
对任意 $\lambda\in[0,1]$，取二进有理数列 $\lambda_j\to\lambda$。由 $f$ 的连续性令 $j\to\infty$，即得
\[
  f(\lambda x+(1-\lambda)y)
  \le\lambda f(x)+(1-\lambda)f(y).
\]
这正是 §8.4 的下凸定义。
\end{xsolution}

\paragraph{第二组参考题 9}
\begin{xsolution}
只需证明每一点都连续。固定 $c\in(a,b)$。题设保证左右极限
\[
  L_-=f(c-),\qquad L_+=f(c+)
\]
存在且有限。Jensen 不等式用于 $c-h,c+h$，令 $h\to0+$，得
\[
  f(c)\le\frac{L_-+L_+}{2}.
\tag{1}
\]
再用于 $c$ 与 $c+2h$，有
\[
  f(c+h)\le\frac{f(c)+f(c+2h)}2.
\]
令 $h\to0+$，得到
\[
  L_+\le\frac{f(c)+L_+}{2},
  \qquad L_+\le f(c).
\]
同理 $L_-\le f(c)$。代入 (1) 即知
\[
  f(c)\le\frac{L_-+L_+}{2}\le f(c),
\]
故
\[
  L_-=L_+=f(c).
\]
所以 $f$ 在 $c$ 连续。$c$ 任意，故 $f$ 在 $(a,b)$ 上连续。
\end{xsolution}

\paragraph{第二组参考题 10}
\begin{xsolution}
固定 $c\in(a,b)$，取 $r>0$ 使 $[c-r,c+r]\subset(a,b)$。由题设存在 $M>0$，使
\[
  |f(x)|\le M\qquad(c-r\le x\le c+r).
\]
由 Jensen 中点不等式反复迭代，对任意整数 $m\ge0$ 有
\[
  f\left((1-2^{-m})c+2^{-m}y\right)
  \le(1-2^{-m})f(c)+2^{-m}f(y).
\tag{1}
\]
给定充分小的 $h\ne0$，选 $m$ 使
\[
  \frac r2<2^m|h|\le r,
\]
并取 $y=c+2^m h$。由 (1)，
\[
  f(c+h)-f(c)
  \le2^{-m}[f(y)-f(c)]
  \le2M\,2^{-m}
  \le\frac{4M}{r}|h|.
\]
对 $-h$ 同样成立。再由
\[
  f(c)\le\frac{f(c+h)+f(c-h)}2
\]
可把上界转成相应下界，从而得到
\[
  |f(c+h)-f(c)|\le C|h|
\]
（$h$ 充分小时），其中 $C$ 与 $c,r,M$ 有关。因此 $f$ 在 $c$ 连续。$c$ 任意，故 $f$ 在 $(a,b)$ 上连续。
\end{xsolution}

\paragraph{第二组参考题 11}
\begin{xsolution}
先在 $[0,1)$ 上证明。令
\[
  H(x)=\int_0^x\bigl(|f(t)|+|f'(t)|\bigr)\,\dd t.
\]
由 $f(0)=f'(0)=0$ 与题设，
\[
  |f'(x)|
  =\left|\int_0^x f''(t)\,\dd t\right|
  \le H(x),
\]
而
\[
  |f(x)|
  =\left|\int_0^x f'(t)\,\dd t\right|
  \le H(x).
\]
故
\[
  H'(x)=|f(x)|+|f'(x)|\le2H(x),
  \qquad H(0)=0.
\]
于是
\[
  (\ee^{-2x}H(x))'\le0.
\]
又 $H\ge0$，所以只能有 $H\equiv0$，从而 $f\equiv0$ 于 $[0,1)$。对 $(-1,0]$ 作同样的反向积分论证，即得
\[
  \boxed{f(x)\equiv0\quad(-1<x<1)}.
\]
\end{xsolution}

\paragraph{第二组参考题 12}
\begin{xsolution}
反设 $f$ 不单调。连续函数在区间上不单调时，可找到 $x_1<x_2<x_3$，使例如
\[
  f(x_2)>\max\{f(x_1),f(x_3)\}.
\]
取
\[
  y\in\bigl(\max\{f(x_1),f(x_3)\},f(x_2)\bigr).
\]
令 $u$ 为 $[x_1,x_2]$ 上第一次达到值 $y$ 的点，$v$ 为 $[x_2,x_3]$ 上最后一次达到值 $y$ 的点。由第一次、最后一次达到的性质，利用单侧差商并结合可微性，有
\[
  f'(u)\ge0,
  \qquad
  f'(v)\le0.
\]
但 $f(u)=f(v)=y$，故微分方程给出
\[
  f'(u)=g(y)=f'(v),
\]
于是二者都等于 0。对上述 $y$ 所在开区间内的每个值重复这一论证，得到 $g(y)=0$。因此当 $f(x)$ 的值落在该开区间内时均有 $f'(x)=0$，由中值定理 $f$ 不可能从端点值上升到 $f(x_2)$，矛盾。故 $f$ 必为单调函数。
\end{xsolution}

\paragraph{第二组参考题 13}
\begin{xsolution}
反设三式对所有 $x\in\RR$ 同时成立。固定 $x_0$。由于 $f''<0$，$f'$ 严格减少；又 $f'(x_0)>0$，所以对每个 $x<x_0$，
\[
  f'(x)>f'(x_0)>0.
\]
于是
\[
  f(x_0)-f(x)
  =\int_x^{x_0}f'(t)\,\dd t
  >(x_0-x)f'(x_0),
\]
即
\[
  f(x)<f(x_0)-(x_0-x)f'(x_0)\to-\infty
  \qquad(x\to-\infty),
\]
与 $f(x)>0$ 矛盾。
\end{xsolution}

\paragraph{第二组参考题 14}
\begin{xsolution}
反设对每个 $x$ 都有
\[
  f(x)f'(x)f''(x)f'''(x)<0.
\]
则四个因子都处处非零，并由连续性或导数的 Darboux 性质各自保持定号。把 $f$ 换成 $-f$ 不改变乘积符号，把 $x$ 换成 $-x$ 也不改变该乘积符号，因此不妨设
\[
  f>0,\qquad f'>0.
\]
于是 $f''$ 与 $f'''$ 异号。

若 $f''>0$、$f'''<0$，则 $f''$ 严格减少。对 $x<0$ 有 $f''(x)>f''(0)>0$，所以
\[
  f'(x)<f'(0)+x f''(0),
\]
右端在 $x\to-\infty$ 时趋于 $-\infty$，与 $f'>0$ 矛盾。

若 $f''<0$、$f'''>0$，则 $f'$ 严格减少。对 $x<0$ 有 $f'(x)>f'(0)>0$，从而
\[
  f(x)<f(0)+x f'(0)\to-\infty
  \qquad(x\to-\infty),
\]
与 $f>0$ 矛盾。

故反设不成立，必存在 $a\in\RR$ 使
\[
  \boxed{f(a)f'(a)f''(a)f'''(a)\ge0}.
\]
\end{xsolution}

\paragraph{第二组参考题 15}
\begin{xsolution}
记微分算子为 $D$。有因式分解
\[
  D^3-D^2-D+1=(D-1)^2(D+1).
\]
先用两个简单引理：若多项式 $h$ 满足 $h'+h\ge0$ 于全实轴，则 $h\ge0$；若 $h'-h\ge0$，则 $h\le0$。例如第一条中，$h'+h$ 与 $h$ 有相同最高次项；若它非零且处处非负，则 $h$ 也有偶数次数和正首项系数，从而有全局最小值。若该最小值为负，在最小点 $h'=0$，便有 $h'+h<0$，矛盾。第二条同理在全局最大点证明。相反号的不等式也相应反向。

令
\[
  q=(D-1)^2(D+1)p\ge0,
  \quad r=(D-1)(D+1)p,
  \quad s=(D+1)p.
\]
由
\[
  r'-r=q\ge0
\]
得 $r\le0$；再由
\[
  s'-s=r\le0
\]
得 $s\ge0$；最后由
\[
  p'+p=s\ge0
\]
得
\[
  \boxed{p(x)\ge0\qquad(\forall x\in\RR)}.
\]
\end{xsolution}

\paragraph{第二组参考题 16}
\begin{xsolution}
首先因式分解
\[
  Q(x)=\bigl(P(x)+xP'(x)\bigr)\bigl(xP(x)+P'(x)\bigr).
\]
将 $P$ 的 $n$ 个大于 1 的互异实根按大小记为
\[
  1<r_1<r_2<\cdots<r_n,
\]
并可取它们为这些根中的相邻根。

第一因子
\[
  P+xP'=(xP)'.
\]
函数 $xP(x)$ 在 $0,r_1,\ldots,r_n$ 处均为 0，因此 Rolle 定理在
\[
  (0,r_1),(r_1,r_2),\ldots,(r_{n-1},r_n)
\]
中各给出至少一个第一因子的零点，共至少 $n$ 个。

第二因子
\[
  xP+P'=\ee^{-x^2/2}\bigl(\ee^{x^2/2}P(x)\bigr)'.
\]
函数 $\ee^{x^2/2}P(x)$ 在 $r_1,\ldots,r_n$ 处为 0，因此在每个 $(r_i,r_{i+1})$ 中至少有一个第二因子的零点，共至少 $n-1$ 个。

还需说明这些零点不会在同一开区间内重合。若某个 $x>1$ 同时满足
\[
  P+xP'=0,\qquad xP+P'=0,
\]
则因系数矩阵行列式 $1-x^2\ne0$，必有 $P(x)=P'(x)=0$，即 $x$ 是 $P$ 的重根。若 $r_i,r_{i+1}$ 取为相邻的互异实根，则区间内不存在这样的根；若原多项式还有额外实根，只需把所有大于 1 的互异根先排序后从中取连续的 $n$ 个即可。因此上述两类零点互异。

故 $Q$ 至少有
\[
  n+(n-1)=\boxed{2n-1}
\]
个互异实根。
\end{xsolution}

\paragraph{第二组参考题 17}
\begin{xsolution}
不动点方程为
\[
  a^x=x.
\]
由于 $a>1$，解必有 $x>1$。取对数得到
\[
  \frac{\ln x}{x}=\ln a.
\]
令
\[
  \phi(x)=\frac{\ln x}{x},\qquad x>1.
\]
则
\[
  \phi'(x)=\frac{1-\ln x}{x^2},
\]
故 $\phi$ 在 $(1,\ee)$ 增加、在 $(\ee,+\infty)$ 减少，唯一最大值为
\[
  \phi(\ee)=\frac1\ee.
\]
因此：若 $\ln a>1/\ee$，无解；若 $\ln a=1/\ee$，恰有一解 $x=\ee$；若 $0<\ln a<1/\ee$，恰有两解。即
\begin{enumerate}[label=(\arabic*)]
  \item $a>\ee^{1/\ee}$ 时无不动点；
  \item $a=\ee^{1/\ee}$ 时恰有一个不动点；
  \item $1<a<\ee^{1/\ee}$ 时有两个不动点。
\end{enumerate}
\end{xsolution}

\paragraph{第二组参考题 18}
\begin{xsolution}
令
\[
  b=-\ln a>0.
\]
则
\[
  g(x)=a^{a^x}=\exp(-b\ee^{-bx}).
\]
不动点必须满足 $x>0$，取对数后等价于
\[
  F(x):=\ln x+b\ee^{-bx}=0.
\]
其导数为
\[
  F'(x)=\frac{1-b^2x\ee^{-bx}}x.
\]
令 $y=bx>0$，则临界点条件为
\[
  y\ee^{-y}=\frac1b.
\]
函数 $y\ee^{-y}$ 在 $y=1$ 处取唯一最大值 $1/\ee$。

若 $b<\ee$，则 $F'(x)>0$ 对所有 $x>0$ 成立；又 $F(0+)=-\infty$、$F(+\infty)=+\infty$，故只有一个零点。$b=\ee$ 时导数仅在 $y=1$ 处为 0，$F$ 仍单调增加，仍只有一个零点。

若 $b>\ee$，则有两个临界参数
\[
  0<y_1<1<y_2.
\]
在临界点，利用 $b=e^y/y$ 得
\[
  F\left(\frac yb\right)
  =2\ln y-y+\frac1y=:H(y).
\]
而
\[
  H'(y)=-\frac{(y-1)^2}{y^2}\le0,
  \qquad H(1)=0.
\]
故 $H(y_1)>0$、$H(y_2)<0$。于是 $F$ 从 $-\infty$ 上升到正的极大值，再下降到负的极小值，最后上升到 $+\infty$，恰有三个零点。

恢复 $b=-\ln a$：
\[
  b\le\ee\iff a\ge\ee^{-\ee}.
\]
因此
\begin{enumerate}[label=(\arabic*)]
  \item $\ee^{-\ee}\le a<1$ 时 $g$ 只有一个不动点；
  \item $0<a<\ee^{-\ee}$ 时 $g$ 恰有三个不动点。
\end{enumerate}
\end{xsolution}

\paragraph{第二组参考题 19}
\begin{xsolution}
记
\[
  F(\theta)=a\sin\theta+b\cos\theta-\sin\theta\cos\theta.
\]
根数只有在出现重根时才可能改变。重根满足
\[
  F(\theta)=0,
  \qquad
  F'(\theta)=a\cos\theta-b\sin\theta-\cos2\theta=0.
\]
联立解出参数
\[
  a=\cos^3\theta,
  \qquad
  b=\sin^3\theta.
\]
故参数平面中的判别曲线是星形线
\[
  |a|^{2/3}+|b|^{2/3}=1.
\]
在这条曲线之外所有根均为单根，因此根数在每个连通区域内保持不变。

在星形线内部取 $(a,b)=(0,0)$，方程化为
\[
  \sin\theta\cos\theta=0,
\]
在 $[0,2\pi)$ 有 4 个根，所以内部处处有 4 个互异实根。

在外部取 $(a,b)=(2,0)$，方程为
\[
  \sin\theta(2-\cos\theta)=0,
\]
只有 $0,\pi$ 两根，所以外部处处有 2 个互异实根。

在星形线的光滑部分，即 $ab\ne0$，重根处
\[
  F''(\theta)=3\sin\theta\cos\theta\ne0,
\]
故恰有一个二重根，其余两根保持为单根，于是共有 3 个互异实根。四个尖点为
\[
  (\pm1,0),\qquad(0,\pm1).
\]
例如在 $(1,0)$ 时
\[
  F(\theta)=\sin\theta(1-\cos\theta),
\]
只有 $0,\pi$ 两个互异根，其中 $0$ 为三重根；其余尖点同理。

综上，令
\[
  \rho=|a|^{2/3}+|b|^{2/3},
\]
则互异实根数为
\[
  \boxed{
  \begin{cases}
    4,&\rho<1,\\
    3,&\rho=1,\ ab\ne0,\\
    2,&\rho>1,\\
    2,&\rho=1,\ ab=0.
  \end{cases}}
\]
\end{xsolution}

\paragraph{第二组参考题 20}
\begin{xsolution}
设椭圆为
\[
  \frac{x^2}{A^2}+\frac{y^2}{B^2}=1,
  \qquad A>B>0,
\]
定点为 $(u,v)$，椭圆上的接触点参数化为
\[
  (A\cos\theta,B\sin\theta).
\]
该点的切向量为 $(-A\sin\theta,B\cos\theta)$。从 $(u,v)$ 到接触点的连线为法线，当且仅当它与切向量垂直，即
\[
  (u-A\cos\theta)(-A\sin\theta)
  +(v-B\sin\theta)(B\cos\theta)=0.
\]
令 $d=A^2-B^2>0$，整理为
\[
  \frac{Au}{d}\sin\theta
  -\frac{Bv}{d}\cos\theta
  -\sin\theta\cos\theta=0.
\]
这正是上一题，其中
\[
  a=\frac{Au}{d},
  \qquad
  b=-\frac{Bv}{d}.
\]
因此判别曲线为
\[
  \left|\frac{u}{d/A}\right|^{2/3}
  +\left|\frac{v}{d/B}\right|^{2/3}=1,
\]
也就是椭圆的渐屈线（星形线）
\[
  u=\frac dA\cos^3\theta,
  \qquad
  v=-\frac dB\sin^3\theta.
\]

所以，从定点到椭圆的法线接触点个数为：渐屈线内部 4 个；外部 2 个；光滑边界上 3 个；四个尖点上 2 个互异接触点（其中一个对应高重数法线）。因此
\[
  \boxed{\text{法线条数最多为 }4,\text{ 且在椭圆渐屈线的内部区域达到。}}
\]
这里“条数”按不同接触点所确定的法线计数。
\end{xsolution}
